\section*{柱体-棱柱和圆柱}

\begin{enumerate}

\item 已知 M= \{正四棱柱\}，N=\{长方体\}，P=\{直四棱柱\}，Q=\{正方体\}，则下列关系中正确的是 \dotfill{(\kh{B})}

\xx{${Q} \supset  {M} \supset  {N} \supset  {{P}}$ }
{${Q} \subset  {M} \subset  {N} \subset  {P}$}
{${Q} \supset  {N} \supset  {M} \supset  {P}$} 
{${Q} \subset  {N} \subset  {M} \subset  {P}$}

\item 下列命题中，正确的是\dotfill{(\kh{C})}

\xx{有一个侧面是矩形的棱柱是直棱柱;}
{有两个侧面是矩形的棱柱是直棱柱;}
{有相邻两个侧面是矩形的棱柱是直棱柱;}
{底面是正多边形的棱柱是直棱柱.}

\item 下列几何体中，哪个是长方体 \dotfill{(\kh{D})}

\xx{直平行六面体} 
{侧面是矩形的四棱柱}
{对角面是全等矩形的四棱柱}
{底面是矩形的直棱柱}

\item 一个棱柱为正四棱柱的条件是 \dotfill{(\kh{C})}

\xx{底面是正方形，有两个侧面垂直于底面}
{底面是正方形，有两个侧面是矩形}
{底面是菱形且有一个顶点处的三条棱两两垂直} 
{每个底面是全等的矩形}

\item 在正方体 ${{ABCD}} - {{A}}_{1}{{B}}_{1}{{C}}_{1}{{D}}_{1}$ 中，与正方体的一个面上的一条对角线成 ${60}^{ \circ  }$ 角的异面的面对角线共有 \dotfill{(\kh{B})}

\xx{2 条} 
{4 条} 
{6 条} 
{8 条}

\item 若长方体的一条对角线与一个顶点处的三条棱成的角分别是 $a,\beta ,\gamma$ ，则有 \dotfill{(\kh{C})}

\xx{${\cos }^{2}a + {\cos }^{2}\beta  + {\cos }^{2}\gamma  = 2$}
{${\sin }^{2}a + {\sin }^{2}\beta  + {\sin }^{2}\gamma  = 3$}
{${\cos }^{2}a + {\cos }^{2}\beta  + {\cos }^{2}\gamma  = 1$}
{${\sin }^{2}a + {\sin }^{2}\beta  + {\sin }^{2}\gamma  = 1$}

\item 如图， $P$ 是平行四边形 ${ABCD}$ 所在平面外一点， $E$ 、 $F$ 分别是 ${PC}$ 、 ${PD}$ 的中点，已知 ${PD}\bot {CD}$ ，且 ${PD} = {2CD}$ ，则异面直线 ${PC}$ 与 ${AB}$ 所成角的大小为\blkda{$arctan2$}.

\begin{figure}[h]
    \centering
    \includegraphics[width=0.25\linewidth]{bodyxb//src/shgoauherfgahbrfbhar.png}
\end{figure}

\item 立体几何中，用一个平面去截一个几何体得到的平面图形叫截面. 已知正方体 ${ABCD} - {A}_{1}{B}_{1}{C}_{1}{D}_{1}$ 的内 切球 $O$ 的直径为 2,过球 $O$ 的一条直径作该正方体的截面,所得的截面面积的最大值为\dotfill{(\kh{D})}

\xx{2} {4 } {$3\sqrt{3}$} { $4\sqrt{2}$}



\item 如图，棱长为 2 的正方体 ${ABCD} - {A}_{1}{B}_{1}{C}_{1}{D}_{1}$ 中， $E$ 为 $C{C}_{1}$ 的中点， 点 $P\text{ 、 }Q$ 分别为面 ${A}_{1}{B}_{1}{C}_{1}{D}_{1}$ 和线段 ${B}_{1}C$ 上动点，则 $\bigtriangleup  {PEQ}$ 周长的最小值为 \dotfill{(\kh{B})}

\begin{figure}[h]
    \centering
    \includegraphics[width=0.25\linewidth]{bodyxb//src/falhr'oaqb'ewihbfqw.png}
\end{figure}

\xx{$2\sqrt{2}$} {$\sqrt{10}$} {$\sqrt{11}$} {$\sqrt{12}$}



\item 正三棱柱 ${ABC} - {A}_{1}{B}_{1}{C}_{1}$ 的底面边长为 8,面对角线 $B{C}_{1} = {10}$ ，D 为 AC 的中点.
\begin{enumerate}
    \item 求证： $A{B}_{1}//$ 平面 ${C}_{1}{BD}$； 
    \item 求异面直线 $A{B}_{1}$ 与 $B{C}_{1}$ 所成的角.
\end{enumerate}

\begin{jieda}
    \begin{enumerate}
        \item 证明：连 $ {B}_{1}{C} $交 $ B{C}_{1} $于 $ {E} $,连 $ {{DE}} $,则 $ {{DE}} $।II $ A{B}_{1} $,
而 $ {{DE}} \subset $面 $ {{C}}_{1}{{DB}},A{B}_{1} ⊄ $面 $ {{C}}_{1}{{DB}},\;\therefore A{B}_{1}// $平面 $ {C}_{1}{DB} $.
        \item 解：由(1)知 $ \angle {DEB} $为异面直线 $ A{B}_{1} $与 $ B{C}_{1} $所成的角,
在 $ {\Delta DEB} $中, $ {DE} = 5,{BD} = 4\sqrt{3},{BE} = 5,\cos \angle {DEB} = \dfrac{{50} - {48}}{2 \times  5 \times  5} = \dfrac{1}{25} $,所以所成角为$arccos\dfrac{1}{25}$.
    \end{enumerate}
\end{jieda}

    \item 如图，已知点 $P$ 在圆柱 $O{O}_{1}$ 的底面圆 $O$ 上， $\angle {AOP} = {120}^{ \circ  }$ ，圆 $O$ 的直径 ${AB} = 4$ ，圆柱的高 $O{O}_{1} = 3$ .
    \begin{enumerate}
        \item 求圆柱的轴截面的面积；
        \item 求点 $A$ 到平面 ${A}_{1}{PO}$ 的距离.
    \end{enumerate}
        
    \begin{figure}[h]
        \centering
        \includegraphics[width=0.25\linewidth]{bodyxb//src/gsljghsejrnfbjhebrfq.png}
    \end{figure}

    \begin{jieda}
        \begin{enumerate}
            \item 面积为12
            \item 以 $ O $为坐标原点, $ {AB} $垂直平分线为 $ x $轴, $ {AB} $为 $ y $轴, $ O{O}_{1} $为 $ z $轴建立空间直角坐标系,
                则 $ \vv{O{A}_{1}} = \left( {0, - 2,3}\right) ,\vv{OP} = \left( {\sqrt{3},1,0}\right) $,
                设平面 $ {A}_{1}{PO} $的一个法向量为 $ \vv{n} = \left( {a,b,c}\right) $,
                则有: $ \left\{  \begin{array}{l} \vv{n} \cdot  {\vv{OA}}_{1} = 0 \\  \vv{n} \cdot  \vv{OP} = 0 \end{array}\right.  \Rightarrow  \left\{  \begin{array}{l}  - {2b} + {3c} = 0 \\  \sqrt{3}a + b = 0 \end{array}\right. $.
                取 $ b = 3 $,
                $ \therefore \vv{n} = \left( {-\sqrt{3},3,2}\right) $.
                则有: $ \vv{A{A}_{1}} \cdot  \vv{n} = \left| \vv{A{A}_{1}}\right|  \cdot  \left| \vv{n}\right|  \cdot  \cos \theta $.
                其中点 $ A $到平面 $ {A}_{1}{PO} $的距离为 $ h = \dfrac{\left| \vv{A{A}_{1}} \cdot  \vv{n}\right| }{\left| \vv{n}\right| } = \dfrac{\left| \left( 0,0,3\right)  \cdot  \left( -\sqrt{3},3,2\right) \right| }{\sqrt{3 + 9 + 4}} = \dfrac{3}{2} $.
                    \end{enumerate}
    \end{jieda}


    
\end{enumerate}

【立体几何综合题】

\begin{enumerate}

\item 已知正方体 ${ABCD} - {A}^{\prime }{B}^{\prime }{C}^{\prime }{D}^{\prime }$ . 求:

\begin{enumerate}
    \item 直线 ${AC}$ 与 ${B}^{\prime }{D}^{\prime }$ 所成的角的大小；
    \item 直线 ${B}^{\prime }D$ 与平面 ${AC}{D}^{\prime }$ 所成的角的大小；
    \item 二面角 $A - C{D}^{\prime } - {B}^{\prime }$ 的大小.
\end{enumerate}

\begin{jieda}
    \begin{enumerate}
        \item $\dfrac{\pi}{2}$；        
        \item $\dfrac{\pi}{2}$；
        \item $arccos\dfrac{1}{3}$
    \end{enumerate}
\end{jieda}

\item 在直三棱柱 ${ABC} - {A}_{1}{B}_{1}{C}_{1}$ 中，底面是等腰直角三角形， $\angle {ACB} = \dfrac{\pi }{2}$ ，侧棱 $A{A}_{1} = 2,D\text{ 、 }E$ 分别是 $C{C}_{1}$ 与 ${A}_{1}B$ 的中点，点 $E$ 在平面 ${ABD}$ 上的射影是 $\bigtriangleup  {ABD}$ 的重心 $G$ . 求直线 ${A}_{1}B$ 与平面 ${ABD}$ 所成角的大小.

\begin{jieda}
    \begin{figure}[h]
        \centering
        \includegraphics[width=0.25\linewidth]{bodyxb//src/flafajlbf'kahef;kbhaw.png}
    \end{figure}

    设 $ \left| {AC}\right|  = \left| {BC}\right|  = a > 0 $,
则 $ A\left( {a,0,0}\right) ,B\left( {0,a,0}\right) ,D\left( {0,0,1}\right) ,G\left( {\dfrac{a}{3},\dfrac{a}{3},\dfrac{1}{3}}\right) ,{A}_{1}\left( {a,0,2}\right) ,E\left( {\dfrac{a}{2},\dfrac{a}{2},1}\right) $,
$ \therefore \vv{EG} = \left( {-\dfrac{a}{6}, - \dfrac{a}{6}, - \dfrac{2}{3}}\right) ,\vv{AD} = \left( {-a,0,1}\right) $,
$ \because {EG} \bot $平面 $ {ABD}，\therefore {EG} \bot  {AD},\therefore \vv{EG} \cdot  \vv{AD} = \dfrac{{a}^{2}}{6} - \dfrac{2}{3} = 0 $,解得： $ a = 2 $.
(1) $ \because {A}_{1}B $在平面 $ {ABD} $上的射影为 $ {BG} $, $ \therefore {A}_{1}B $与平面 $ {ABD} $所成角为 $ \angle {A}_{1}{BG} $.
$ \because \vv{B{A}_{1}} = \left( {2, - 2,2}\right) ,\vv{BG} = \left( {\dfrac{2}{3}, - \dfrac{4}{3},\dfrac{1}{3}}\right) $,
$ \therefore \cos \angle {A}_{1}{BG} = \dfrac{\left| \vv{B{A}_{1}} \cdot  \vv{BG}\right| }{\left| \vv{B{A}_{1}}\right|  \cdot  \left| \vv{BG}\right| } = \dfrac{\left| \dfrac{4}{3} + \dfrac{8}{3} + \dfrac{2}{3}\right| }{2\sqrt{3} \times  \dfrac{\sqrt{21}}{3}} = \dfrac{\sqrt{7}}{3} $,
$ \therefore {A}_{1}B $与平面 $ {ABD} $所成角大小为 $ \arccos \dfrac{\sqrt{7}}{3} $.
\end{jieda}



\item 已知正方体 ${ABCD}$ - ${PQMN}$ 的棱长为 $a$ ,求:
\begin{enumerate}
    \item 异面直线 ${MD}$ 与 ${PC}$ 的距离；
    \item 异面直线 ${MD}$ 与 ${QC}$ 的距离.
\end{enumerate}

\begin{jieda}
    可用建系的方法求解两个小题
    \begin{enumerate}
        \item $\dfrac{\sqrt{6}a}{6}$
        \item $\dfrac{\sqrt{3}a}{3}$
    \end{enumerate}
\end{jieda}

\item 如图，平面 ${{PAD}} \bot$ 平面 ${{ABCD}}$ , ${{ABCD}}$ 为正方形， $\angle {{PAD}} = {90}^{ \circ  }$ ,且 ${{PA}} = {{AD}} = 2,{\;E}\text{ 、 }{\;F}$分别是线段 PA、PD、CD 的中点.

\begin{figure}[h]
    \centering
    \includegraphics[width=0.25\linewidth]{bodyxb//src/aofhoawhofbqwohbrfq.png}
\end{figure}

\begin{enumerate}
    \item 求证: PB//平面EFG;
    \item 求异面直线 EG 与 BD 所成的角的余弦值；
    \item 在线段 CD 上是否存在一点 Q，使得点 A 到平面 EFQ 的距离为 $\dfrac{4}{5}$ ，若存在，求出 CQ 的值；若不存在，请说明理由.
\end{enumerate}

\begin{jieda}
    \begin{enumerate}
        \item 取 $ {AB} $中点 $ H $,连接 $ {GH} $, $ {HE} $,
$ \because E\text{、}F\text{、}G $分别是线段 $ {PA}\text{、}{PD}\text{、}{CD} $的中点,
$ \therefore {GH}//{AD}//{EF} $.
$ \therefore E,F,G,H $四点共面.
又 $ H $为 $ {AB} $中点,
$ \therefore {EH}//{PB} $.
又 $ {EH} \subset $平面 $ {EFG},{PB} ⊄ $平面 $ {EFG} $,
$ \therefore {PB}// $平面 $ {EFG} $.
        \begin{figure}[h]
            \centering
            \includegraphics[width=0.25\linewidth]{bodyxb//src/flejrhfoqbrfbqejrf.png}
        \end{figure}


        \item 因为平面 $ {PAD} \bot $平面 $ {ABCD} $, $ \angle {PAD} = {90}^{ \circ  } $,
即 $ {PA}\bot {AD} $,又平面 $ {PAD} \cap $平面 $ {ABCD} = {AD} $, $ {PA} \subset $平面 $ {PAD} $,所以 $ {PA}\bot $平面 $ {ABCD} $,又 $ {AB} \subset $平面 $ {ABCD} $,所以 $ {PA}\bot {AB} $,
又 $ {ABCD} $为正方形,即 $ {AB}\bot {AD} $,
如图建立空间直角坐标系,

        \begin{figure}[h]
            \centering
            \includegraphics[width=0.25\linewidth]{bodyxb//src/pojfaljwnfoajnjwofa.png}
        \end{figure}

        则 $ E\left( {0,0,1}\right) ,A\left( {0,0,0}\right) ,B\left( {2,0,0}\right) ,D\left( {0,2,0}\right) ,G\left( {1,2,0}\right) ,F\left( {0,1,1}\right) $,
所以 $ \vv{EG} = \left( {1,2, - 1}\right) ,\vv{BD} = \left( {-2,2,0}\right) $,
所以 $ \cos \langle \vv{EG},\vv{BD}\rangle  = \dfrac{\vv{EG} \cdot  \vv{BD}}{\left| \vv{EG}\right| \left| \vv{BD}\right| } = \dfrac{-2 + 4}{\sqrt{6} \times  2\sqrt{2}} = \dfrac{\sqrt{3}}{6} $,
即异面直线 $ {EG} $与 $ {BD} $所成的角的余弦值为 $ \dfrac{\sqrt{3}}{6} $.


    \item 假设在线段 $ {CD} $上存在一点 $ Q $满足题设条件,令 $ {CQ} = m\left( {0 \leq  m \leq  2}\right) $,
则 $ {DQ} = 2 - m $,
$ \therefore {点Q} $的坐标为 $ \left( {2 - m,2,0}\right) ,\vv{EQ} = \left( {2 - m,2, - 1}\right) ,\vv{EF} = \left( {0,1,0}\right) $,
设平面 $ {EFQ} $的法向量 $ \vv{n} = \left( {x,y,z}\right) $,
则 $ \left\{  \begin{array}{l} \vv{n} \cdot  \vv{EF} = y = 0 \\  \vv{n} \cdot  \vv{EQ} = \left( {2 - m}\right) x + {2y} - z = 0 \end{array}\right. $,取 $ \vv{n} = \left( {1,0,2 - m}\right) $,
又 $ \vv{AE} = \left( {0,0,1}\right) $,
所以点 $ A $到平面 $ {EFQ} $的距离 $ d = \dfrac{\left| \vv{AE} \cdot  \vv{n}\right| }{\left| \vv{n}\right| } = \dfrac{\left| 2 - m\right| }{\sqrt{1 + {\left( 2 - m\right) }^{2}}} = \dfrac{4}{5} $,
即 $ {\left( 2 - m\right) }^{2} = \dfrac{16}{9} $,解得 $ m = \dfrac{2}{3} $或 $ m = \dfrac{10}{3} $(舍去),
故存在点 $ Q $,当 $ {CQ} = \dfrac{2}{3} $时,点 $ A $到平面 $ {EFQ} $的距离为 $ \dfrac{4}{5} $.

    \end{enumerate}
\end{jieda}


\end{enumerate}





\section*{柱体的体积和表面积}

A组：
\begin{enumerate}
    \item 有三个命题:甲:底面是平行四边形的四棱柱是平行六面体；乙:底面是矩形是平行六面体是长方体; 丙: 直四棱柱是直平行六面体; 其中真命题的个数是\dotfill{(\kh{B})}
    \xx{0 个}{1 个}{2 个}{3 个}
    \begin{jieda}
        甲：正确，乙反例：斜棱柱，丙反例：底面是梯形
    \end{jieda}

    \item 若一个长方体其一顶点的三个面的面积分别为 $\sqrt{2}$ 、 $\sqrt{3}$ 、 $\sqrt{6}$ ，则这个长方体的对角线的长是\blkda{$\sqrt{6}$}.
    \begin{jieda}
        棱长：$1, \sqrt{2}, \sqrt{3}$
    \end{jieda}

    \item 已知长方体的全面积是 11，十二条棱长度之和是 24，则这个长方体的一条对角线长\blkda{5}
    \begin{jieda}
        设棱长分别为$a, b, c$，则$\left\{\begin{array}{l}
             2(ab+bc+ac) = 11  \\
             4(a+b+c) = 24
        \end{array}\right.$，$a^2+b^2+c^2 = (a+b+c)^2-2(ab+bc+ac) = 36-11 = 25$，故对角线长度为$5$.
    \end{jieda}

    \item 下列各图中不是正方体表面展开图的是\dotfill{(\kh{C})}
    \xx{\includegraphics[width=0.2\textwidth]{bodyxb/src/bo_d5aaqaf7aajc73815vqg_0_310_920_143_92_0.jpg}}{\includegraphics[width=0.2\textwidth]{bodyxb/src/bo_d5aaqaf7aajc73815vqg_0_547_919_137_91_0.jpg}}{\includegraphics[width=0.2\textwidth]{bodyxb/src/bo_d5aaqaf7aajc73815vqg_0_782_915_138_96_0.jpg}}{\includegraphics[width=0.2\textwidth]{bodyxb/src/bo_d5aaqaf7aajc73815vqg_0_1030_945_195_69_0.jpg}}
\end{enumerate}

B组：

\begin{enumerate}
    \item 
    \begin{dingautolist}{172}
        \item 若长方体面上矩形的对角线长分别为 $a$ 、 $b$ 、 $c$ ，则它的体对角线为\blkda[2]{$\sqrt{\dfrac{a^2+b^2+c^2}{2}}$}
        \item 若长方体的全面积为 ${24}{{{cm}}}^{2}$ ，所有棱长之和为 ${24}{{cm}}$ ，则其体对角线长为\blkda{$2\sqrt{3}cm$}
        \item 若长方体的三个面的面积分别是 $\sqrt{2}{{{cm}}}^{2}$ ， $\sqrt{3}{{{cm}}}^{2}$ ， $\sqrt{6}{{{cm}}}^{2}$ ，则其体对角线为\blkda{$\sqrt{6}cm$}.
    \end{dingautolist}
    \item 若长方体的高为 $h$ ，底面积为 ${P}$ ，垂直于底面的对角面的面积为 ${Q}$ ，则此长方体的侧面积是\blkda[3]{$2\sqrt{Q^2+2Ph^2}$}.
    \item 底面是菱形的直棱柱，若对角面的面积为 ${50}\sqrt{2}c{m}^{2}$ 和 ${10}\sqrt{14}c{m}^{2}$ ，而底边长为 ${8cm}$ ， 则其对角线的长等于\blkda[3]{$15cm, 9cm$}.
    \begin{jieda}
        设底面菱形对角线分别为$a, b$，高为$h$（单位$cm$），则有$\left\{\begin{array}{l}
             ah = 50\sqrt{2} \\
             bh = 10\sqrt{14} \\
             a^2+b^2 = 256
        \end{array}\right.$，解得$\left\{\begin{array}{l}
             a = 10\sqrt{2} \\
             b = 2\sqrt{14} \\
             h = 5
        \end{array}\right.$，故对角线长度分别为$15cm, 9cm$.
    \end{jieda}
    \item 已知三个平面两两垂直，且交于同一点 ${P}$ ，空间一点 ${Q}$ 到三个平面的距离分别为 3，4，12，那么 P、Q 两点的距离等于\blkda{$13$}.
    \item 在斜三棱柱 ${{ABC}} - {{A}}_{1}{{B}}_{1}{{C}}_{1}$ 中， $\angle {{BAC}} = {90}^{ \circ  },{{{BC}}}_{1} \bot  {{AC}}$ ，则 ${{C}}_{1}$ 在底面 ${{ABC}}$ 上的射影 ${H}$ 必在\dotfill{(\kh{A})}
    \xx{直线 ${{AB}}$ 上}{直线 $BC$ 上}{直线 $AC$ 上}{$\triangle ABC$ 内部}
    \begin{jieda}
        根据三垂线定理可知，$BC_1$的射影$BH$应与$AB$重合.
    \end{jieda}
    \item 如图，长方体 ${ABCD} - {A}_{1}{B}_{1}{C}_{1}{D}_{1}$ 中， ${AB} = {BC} = 1$ ， ${A{A}_{1}} = 2$ ， $M$ 为 $A{A}_{1}$ 的中点，过 ${B}_{1}M$ 作长方体的截面 $\alpha$ 交棱 $C{C}_{1}$ 于 $N$ ,则真命题个数为\blkda{2}个.
    \begin{center}
    \includegraphics[width=0.2\textwidth]{bodyxb/src/bo_d5aaqaf7aajc73815vqg_1_1047_587_308_184_0.jpg}
    \end{center}
    \begin{dingautolist}{172}
        \item 截面 $\alpha$ 可能为六边形
        \item 存在点 $N$ ,使得 ${BN} \bot$ 截面 $\alpha$
        \item 截面 $\alpha$ 为平行四边形,则 $1 \leq  {CN} \leq  2$
        \item 当 $N$ 与 $C$ 重合时,截面面积为 $\dfrac{3\sqrt{6}}{4}$
    \end{dingautolist}
    \begin{jieda}
        \begin{dingautolist}{172}
            \item 如图可知：

            当 $1 \leq  \left| {CN}\right|  \leq  2$ 时,截面 $\alpha$ 为平行四边形,
            
            当 $0 < {CN} < 1$ 时，截面 $\alpha$ 为五边形，
            
            当 ${CN} = 0$ ，即点 $N$ 与点 $C$ 重合时，截面 $\alpha$ 为梯形

            \begin{center}
            \includegraphics[width=0.3\textwidth]{bodyxb/src/bo_d5achqbef24c73bijjqg_0_243_1759_367_253_0.jpg}
            \end{center}
            
            \ps $B_1$在顶点，截面至多5边形.
            \item 假设存在，则$BN \bot B_1M$，则应有$N$和$C$重合，此时不满足要求，故不存在.
            \item 正确
            \item 因为当 $N$ 与 $C$ 重合时,截面 $\alpha$ 为梯形，如图所示，过 $M$ 作 $M{M}^{\prime }$ 垂直于 ${B}_{1}C$ 于点 ${M}^{\prime }$ ，
            
            \begin{center}
            \includegraphics[width=0.2\textwidth]{bodyxb/src/bo_d5achqbef24c73bijjqg_1_244_901_336_287_0.jpg}
            \end{center}
            
            设梯形的高为 $h,{B}_{1}{M}^{\prime } = x$ ,
            
            则由平面几何知识可得 ${h}^{2} = {\left( \sqrt{2}\right) }^{2} - {x}^{2} = {\left( \dfrac{\sqrt{5}}{2}\right) }^{2} - {\left( \dfrac{\sqrt{5}}{2} - x\right) }^{2}$，解得 $x = \dfrac{2\sqrt{5}}{5},h = \dfrac{\sqrt{30}}{5}$
            
            所以截面 $\alpha$ 的面积为 $\dfrac{\sqrt{5} + \dfrac{\sqrt{5}}{2}}{2} \times  \dfrac{\sqrt{30}}{5} = \dfrac{3\sqrt{6}}{4}$ 
        \end{dingautolist}
    \end{jieda}
    \item 如图，在棱长为 1 的正方体 ${ABCD} - {A}_{1}{B}_{1}{C}_{1}{D}_{1}$ 中， $P$ 为底面 ${ABCD}$ 内(包括边界)的动点，满足 ${D}_{1}P$ 与直线 $C{C}_{1}$ 所成角的大小为 $\dfrac{\pi }{6}$ ， 则线段 ${DP}$ 扫过的面积为\blkda{$\dfrac{\pi}{12}$}.
    \begin{center}
    \includegraphics[width=0.2\textwidth]{bodyxb/src/bo_d5aaqaf7aajc73815vqg_1_231_922_317_286_0.jpg}
    \end{center}
    \begin{jieda}
        扫过的区域为：以$D$为圆心，$\dfrac{\sqrt{3}}{3}$为半径，圆心角为$\dfrac{\pi}{2}$的扇形.
    \end{jieda}

    \item 在底面边长与侧棱长均为 $a$ 的正三棱柱 ${{ABC}} - {{A}}_{1}{{B}}_{1}{{C}}_{1}$ 中，已知 $M$ 为 ${A}_{1}{B}_{1}$ 的中点，则 $M$ 到 ${BC}$ 的距离是\blkda{$\dfrac{\sqrt{19}}{4}a$}.

    \item 已知正三棱柱 ${ABC} - {A}_{1}{B}_{1}{C}_{1}$ 的底面边长为 $2{{cm}}$ ，高为 $4{{cm}}$ ，过 ${BC}$ 作一截面，截面与底面 $ABC$ 成 ${60}^{ \circ  }$ 角，则截面面积是\blkda{$2\sqrt{3} cm^2$}.
    \begin{jieda}
        $cos\ang{60} = \dfrac{S}{S'}$，故$S' = 2S = 2\sqrt{3} cm^2$
    \end{jieda}

    \item 运用祖暅原理计算球体积时，构造一个底面半径和高都与球半径相等的圆柱，与半球( 如图一)放置在同一平面上，然后在圆柱内挖去一个以圆柱下底面圆心为顶点，圆柱上底面为底面的圆锥 (如图二)，用任何一个平行于底面的平面去截它们时，可证得所截得的两个截面面积相等，由此证明该几何体与半球体积相等. (球的体积公式:半径为 ${R}$ 时， $V = \dfrac{4}{3}\pi {R}^{3}$ )现将椭圆 $\dfrac{{x}^{2}}{4} + \dfrac{{y}^{2}}{9} = 1$ 绕 $y$ 轴旋转一周后得一橄榄状的几何体(如图三)，类比上述方法，运用祖暅原理可求得其体积等于\blkda{$16\pi$}.

    \threetu{\includegraphics[width=0.5\textwidth]{bodyxb/src/bo_d5aaqaf7aajc73815vqg_1_329_1931_259_152_0.jpg}}{\includegraphics[width=0.5\textwidth]{bodyxb/src/bo_d5aaqaf7aajc73815vqg_1_630_1913_259_171_0.jpg}}{\includegraphics[width=0.5\textwidth]{bodyxb/src/bo_d5aaqaf7aajc73815vqg_1_919_1847_194_275_0.jpg}}

    \item 在正方体 ${ABCD} - {A}_{1}{B}_{1}{C}_{1}{D}_{1}$ 中， ${AB} = 2,G$ 为棱 ${CD}$ 的中点， $P,Q$ 分别为 $B{C}_{1},C{C}_{1}$ 上的动点，则 ${PQ} + {QG}$ 的最小值为\blkda{$\dfrac{3\sqrt{2}}{2}$}.
    \begin{jieda}
        将正方体的侧面 ${BC}{C}_{1}{B}_{1}$ 与 ${DC}{C}_{1}{D}_{1}$ 展开到同一平面
        
        在同一平面内可知 ${PQ} + {QG}$ 的最小值就是点 $G$ 到 ${B}^{\prime }{C}_{1}$ 的距离,
        
        正方体 ${ABCD} - {A}_{1}{B}_{1}{C}_{1}{D}_{1}$ 中, ${AB} = 2,G$ 为棱 ${CD}$ 的中点,所以 ${CG} = 1,{B}^{\prime }G = 3$ ,
        
        $C{B}^{\prime }{B}_{1}^{\prime }{C}_{1}$ 是正方形,所以 ${GH} = \dfrac{3\sqrt{2}}{2}$
    \end{jieda}
    
    \begin{center}
    \includegraphics[width=0.2\textwidth]{bodyxb/src/bo_d5aaqaf7aajc73815vqg_2_233_216_299_304_0.jpg}
    \end{center}
    
    \item 正三棱柱 ${ABC} - {A}_{1}{B}_{1}{C}_{1}$ 的底边长侧棱长都是 2， $M$ 为 ${AB}$ 的中点， $N$ 为 $C{C}_{1}$ 的中点，则在棱柱表面上，从 $M$ 到 $N$ 的最短路程是\blkda{$\sqrt{4+\sqrt{3}}$}.
    \begin{jieda}
        如图,三棱柱表面由点 $M$ 到 $N$ 的展开图有如下情况
        
        \begin{center}
        \includegraphics[width=0.2\textwidth]{bodyxb/src/bo_d5abfif7aajc738167kg_0_240_576_299_370_0.jpg}
        \end{center}
        
        
        如图,若 ${MN}$ 过 $B{B}_{1}$ 时,此时 ${MN} = \sqrt{{1}^{2} + {3}^{2}} = \sqrt{10}$ ,
        
        \begin{center}
        \includegraphics[width=0.3\textwidth]{bodyxb/src/bo_d5abfif7aajc738167kg_0_242_1048_357_245_0.jpg}
        \end{center}
        
        
        若 ${MN}$ 过 $A{A}_{1}$ 时,与 ${MN}$ 过 $B{B}_{1}$ 一样,此时 ${MN} = \sqrt{{1}^{2} + {3}^{2}} = \sqrt{10}$ ;
        
        第二种情况,当 ${MN}$ 过 ${AC}$ 时, ${NC} = 1,{CM} = \sqrt{3},\angle {NCM} = {90}^{ \circ  } + {30}^{ \circ  } = {120}^{ \circ  }$ ,
        
        ${MN} = \sqrt{N{C}^{2} + C{M}^{2} - {2NC} \cdot  {CM} \cdot  \cos {120}^{ \circ  }} = \sqrt{4 + \sqrt{3}}$
        
        \begin{center}
        \includegraphics[width=0.2\textwidth]{bodyxb/src/bo_d5abfif7aajc738167kg_0_241_1568_236_371_0.jpg}
        \end{center}
        
        
        若 ${MN}$ 过 ${BC}$ 与 ${MN}$ 过 ${AC}$ 一样,此时 ${MN} = \sqrt{4 + \sqrt{3}}$ ,
        
        由 $\sqrt{10} > \sqrt{4 + \sqrt{3}}$
        
        所以从 $M$ 到 $N$ 的最短路程是 $\sqrt{4 + \sqrt{3}}$ .
    \end{jieda}

    \item 已知正方体 ${ABCD} - {A}_{1}{B}_{1}{C}_{1}{D}_{1}$ 的棱长为 2,点 $E$ 是 ${A}_{1}{B}_{1}$ 的中点,则点 $A$ 到直线 ${BE}$ 的距离是\blkda{$\dfrac{4\sqrt{5}}{5}$}.
    \begin{jieda}
        以 $D$ 为原点,以 $\vv{DA},\vv{DC},\vv{D{D}_{1}}$ 的方向为 $x$ 轴、 $y$ 轴、 $z$ 轴的正方向建立空间直角坐标系.
        
        则 $A\left( {2,0,0}\right) ,B\left( {2,2,0}\right) ,E\left( {2,1,2}\right)$ ,
        
        所以 $\vv{BA} = \left( {0, - 2,0}\right) ,\vv{BE} = \left( {0, - 1,2}\right)$ ,
        
        记与 $\vv{BE}$ 同向的单位向量为 $\vv{u}$ ,则 $\vv{u} = \left( {0, - \dfrac{\sqrt{5}}{5},\dfrac{2\sqrt{5}}{5}}\right)$ ,
        
        所以,点 $A$ 到直线 ${BE}$ 的距离 $d = \sqrt{{\vv{BA}}^{2} - {\left( \vv{BA} \cdot  \vv{u}\right) }^{2}} = \sqrt{4 - {\left( \dfrac{2\sqrt{5}}{5}\right) }^{2}} = \dfrac{4\sqrt{5}}{5}$ .

    \end{jieda}


    \item 已知双曲线方程 ${x}^{2} - {y}^{2} = 1$ ，直线 $y = 0$ ， $y = 3$ 在第一象限内与双曲线及渐近线围成的图形绕 $y$ 轴旋转一周所得几何体的体积为\blkda{ ${3\pi }$}. (提示: 利用祖暅原理)
    \begin{jieda}
        双曲线 ${x}^{2} - {y}^{2} = 1$ 的渐近线方程为 $y =  \pm  x$ ,
        
        记 $D$ 绕 $y$ 轴旋转一周所得的几何体为 $\Omega$ ,
        
        $y = 3$ 在第一象限内与渐近线 $y = x$ 的交点 $N$ 的坐标为 $\left( {3,3}\right)$，$y = 3$ 与双曲线 ${x}^{2} - {y}^{2} = 1$ 在第一象限交点 $B$ 的坐标为 $\left( {\sqrt{10},3}\right)$ ,

        在所得几何体中，在高为 $h$ 处作一截面，令 $y = h$ ，则其与渐近线 $y = x$ 的交点坐标为 $\left( {h,h}\right)$ ，
        
        与双曲线的一支 ${x}^{2} - {y}^{2} = 1\left( {x > 0}\right)$ 的交点坐标为 $\left( {\sqrt{1 + {h}^{2}},h}\right)$ ，
        
        则截面面积为 $\pi \left\lbrack  {{\left( \sqrt{1 + {h}^{2}}\right) }^{2} - {h}^{2}}\right\rbrack   = \pi$ ,
        
        利用祖暅原理得 $\Omega$ 的体积相当于底面面积为 $\pi$ 高为 3 的圆柱的体积,
        
        $\therefore \Omega$ 的体积 $V = \pi  \times  3 = {3\pi }$ 

        \begin{center}
        \includegraphics[width=0.3\textwidth]{bodyxb/src/bo_d5abfif7aajc738167kg_5_239_413_399_360_0.jpg}
        \end{center}
    \end{jieda}

    \item 已知三棱柱 ${ABC} - {A}_{1}{B}_{1}{C}_{1}$ 中，所有棱长均为 6，且 $\angle {A}_{1}{AB} = \angle {A}_{1}{AC} = {60}^{ \circ  }$ ，则该三棱柱的侧面积等于\blkda{$36+36\sqrt{3}$}.

    \item 如图代表未折叠正方体的展开图，将其折叠起来，变成正方体后，图形是\dotfill{(\kh{B})}
    
    \begin{center}
    \includegraphics[width=0.2\textwidth]{bodyxb/src/bo_d5aaqaf7aajc73815vqg_2_237_1091_213_171_0.jpg}
    \end{center}

    
    \begin{center}
    \includegraphics[width=0.8\textwidth]{bodyxb/src/bo_d5aaqaf7aajc73815vqg_2_519_1107_528_208_0.jpg}
    \end{center}

    \item 某车间生产一种六角螺母，每一个六角螺母是由棱长和高均为 $2{{cm}}$ 的正六棱柱形的工件加工而成，需要在工件底面的中心处打一个圆柱形通孔(如图)， 若通孔内凹槽忽略不计，则六角螺母表面积的最大值为\dotfill{(\kh{A})}
    \begin{center}
    \includegraphics[width=0.15\textwidth]{bodyxb/src/bo_d5aaqaf7aajc73815vqg_2_241_1331_150_166_0.jpg}
    \end{center}
    \xx{${24} + {12}\sqrt{3} + {2\pi }{{{cm}}}^{2}$}{${24} + {12}\sqrt{3} + {4\pi }{{{cm}}}^{2}$}{${24} + 6\sqrt{3} + {2\pi }{{{cm}}}^{2}$}{${24} + 6\sqrt{3} + {4\pi }{{{cm}}}^{2}$}
    \begin{jieda}
        设通孔的半径为 $r$ ,则六角螺母表面积$S = 2 \times  2 \times  6 + \dfrac{1}{2} \times  2 \times  \sqrt{3} \times  6 \times  2 + {2\pi r} \times  2 - {2\pi }{r}^{2} = {24} + {12}\sqrt{3} - {2\pi }{\left( r - 1\right) }^{2} + {2\pi },$
        因此当 $r = 1$ 时, ${S}_{\max } = {24} + {12}\sqrt{3} + {2\pi }$ ,
        所以六角螺母表面积的最大值为 ${24} + {12}\sqrt{3} + {2\pi }{{{cm}}}^{2}$.
    \end{jieda}

    \item 已知长方体 ${ABCD} - {A}_{1}{B}_{1}{C}_{1}{D}_{1},{AB} = {BC} = 2,B{B}_{1} = 3$ .
            
    \begin{center}
    \includegraphics[width=0.2\textwidth]{bodyxb/src/bo_d5aaqaf7aajc73815vqg_2_231_1720_253_287_0.jpg}
    \end{center}
    \begin{enumerate}
        \item 求二面角 ${B}_{1} - {A}_{1}{C}_{1} - B$ 的大小
        \item 求直线 ${B}_{1}{C}_{1}$ 与平面 ${A}_{1}B{C}_{1}$ 所成角的大小
    \end{enumerate}
    \begin{jieda}
        \begin{enumerate}
            \item 以 $D$ 为原点， ${DA}$ ， ${DC}$ ， $D{D}_{1}$ 分别为 $x$ ， $y$ ， $z$ 轴建系如图
            
            \begin{center}
            \includegraphics[width=0.2\textwidth]{bodyxb/src/bo_d5acvjbef24c73bijrlg_0_244_1622_336_393_0.jpg}
            \end{center}
            
            ${A}_{1}\left( {2,0,3}\right) ,{B}_{1}\left( {2,2,3}\right) ,{C}_{1}\left( {0,2,3}\right) ,B\left( {2,2,0}\right)$
            
            所以 $\vv{{A}_{1}B} = \left( {0,2, - 3}\right) ,\vv{{A}_{1}{C}_{1}} = \left( {-2,2,0}\right)$
            
            设平面 ${A}_{1}B{C}_{1}$ 的一个法向量为 $\vv{{n}_{1}} = \left( {x,y,z}\right)$ ,则
            
            $\left\{  {\begin{array}{l} \vv{{n}_{1}} \cdot  \vv{{A}_{1}B} = 0 \\  \vv{{n}_{1}} \cdot  \vv{{A}_{1}{C}_{1}} = 0 \end{array},\text{ 即 }\left\{  \begin{array}{l} {2y} - {3z} = 0 \\   - {2x} + {2y} = 0 \end{array}\right. }\right.$
            
            取 $x = 1$ ,则 $y = 1,z = \dfrac{2}{3}$
            
            即 $\vv{{n}_{1}} = \left( {1,1,\dfrac{2}{3}}\right)$
            
            显然平面 ${A}_{1}{B}_{1}{C}_{1}$ 的一个法向量为 $\vv{{n}_{2}} = \left( {0,0,1}\right)$
            
            则 $\cos  < {B}_{1} - {A}_{1}{C}_{1} - B >  = \cos  < \vv{{n}_{1}},\vv{{n}_{2}} >  = \dfrac{\vv{{n}_{1}} \cdot  \vv{{n}_{2}}}{\left| \vv{{n}_{1}}\right|  \cdot  \left| \vv{{n}_{2}}\right| } = \dfrac{\dfrac{2}{3}}{\sqrt{1 + 1 + \dfrac{4}{9}} \cdot  1} = \dfrac{\sqrt{22}}{11}$
            
            所以二面角 ${B}_{1} - {A}_{1}{C}_{1} - B$ 的大小为 $\arccos \dfrac{\sqrt{22}}{11}$ .
            
            \item  $\vv{{B}_{1}{C}_{1}} = \left( {-2,0,0}\right)$
            
            则 $\cos  < \vv{{B}_{1}{C}_{1}},\vv{{n}_{1}} >  = \dfrac{\vv{{B}_{1}{C}_{1}} \cdot  \vv{{n}_{1}}}{\left| \vv{{B}_{1}{C}_{1}}\right|  \cdot  \left| \vv{{n}_{1}}\right| } = \dfrac{-2}{2 \cdot  \sqrt{1 + 1 + \dfrac{4}{9}}} =  - \dfrac{3\sqrt{22}}{22}$
            
            记直线 ${B}_{1}{C}_{1}$ 与平面 ${A}_{1}B{C}_{1}$ 所成角为 $\alpha$
            
            则 $\sin \alpha  = \left| {\cos  < \vv{{B}_{1}{C}_{1}},\vv{{n}_{1}} > }\right|  = \dfrac{3\sqrt{22}}{22}$
            
            所以 $\alpha  = \arcsin \dfrac{3\sqrt{22}}{22}$ .
        \end{enumerate}
    \end{jieda}
    
    \item 正三棱柱侧面展开图是边长为 2 和 4 的矩形，则它的表面积是多少？
    \begin{jieda}
        因为正三棱柱的侧面展开图是边长分别为2和4的矩形,所以有以下两种情况：
        \begin{dingautolist}{172}
            \item 当2是下底面的周长,4是正三棱柱的高时,
            正三棱柱的表面积为 ${S}_{\text{表 }} = {S}_{\text{侧 }} + 2{S}_{\text{底 }} = 2 \times  4 + 2 \times  \dfrac{1}{2} \times  {\left( \dfrac{2}{3}\right) }^{2} \times  \dfrac{\sqrt{3}}{2} = 8 + \dfrac{2\sqrt{3}}{9}$ 

            \item 当4是下底面的周长,2是正三棱柱的高时,
            正三棱柱的表面积为 ${S}_{\text{表 }} = {S}_{\text{侧 }} + 2{S}_{\text{底 }} = 4 \times  2 + 2 \times  \dfrac{1}{2} \times  {\left( \dfrac{4}{3}\right) }^{2} \times  \dfrac{\sqrt{3}}{2} = 8 + \dfrac{8\sqrt{3}}{9}$
        \end{dingautolist}
    \end{jieda}

    \item 已知侧面积为 27 的正三棱柱的侧棱恰好是某个圆柱的三条母线，且圆柱底面半径为 2 ，求圆柱的表面积.
    \begin{jieda}
        设正三棱柱的底面边长为 $a$ ,侧棱长为 $h$ ,
        由题意知,正三棱柱是圆柱的内接三棱柱,
        所以正三棱柱的侧棱长 $h$ 即为圆柱的母线长,
        又因为正三棱柱底面三角形是圆柱底面圆的内接三角形,且圆柱底面半径为2,
        所以三角形边长为 $a = 2\sqrt{3}$ ,
        又因为 $S$ 三棱柱侧 $= {3ah} = 3 \times  2\sqrt{3} \times  h = {27}$ ,所以 $h = \dfrac{3\sqrt{3}}{2}$ ,
        所以 ${S}_{\text{圆柱表 }} = {S}_{\text{侧 }} + 2{S}_{\text{底 }} = {2\pi } \times  2 \times  \dfrac{3\sqrt{3}}{2} + 2 \times  \pi  \times  {2}^{2} = \left( {6\sqrt{3} + 8}\right) {\pi }$
    \end{jieda}

    \item 已知平行六面体 ${{ABCD}} - {{A}}_{1}{{B}}_{1}{{C}}_{1}{{D}}_{1}$ 的各棱长都等于 ${a}$ ，且 $\angle {{B}}_{1}{{C}}_{1}{{D}}_{1} = \angle {{{CC}}}_{1}{{B}}_{1} = \angle {{{CC}}}_{1}{{D}}_{1} = \ang{60}$ . 求 ${C}$ 到平面 ${{A}}_{1}{{B}}_{1}{{C}}_{1}$ 的距离.
    \begin{jieda}
        看做边长为$a$的正四面体，顶点到底面的距离，为$\dfrac{\sqrt{6}}{3}a$
    \end{jieda}

    \item 《九章算术》中，将四个面都为直角三角形的四面体称之为鳖臑.在如图所示的空间几何体中,四边形 ${ABCD}$ 为矩形,点 $P$ 为四边形 ${ABCD}$ 所在平面外一点,且 ${PD} \bot$ 平面 ${ABCD},{PD} = {CD}$ , 点 $E$ 是 ${PC}$ 的中点,连接 ${DE}\text{ 、 }{BD}\text{ 、 }{BE}$ .
    
    \begin{center}
    \includegraphics[width=0.2\textwidth]{bodyxb/src/bo_d5aaqaf7aajc73815vqg_3_240_1296_341_314_0.jpg}
    \end{center}
    
    \begin{enumerate}
        \item 证明: ${DE} \bot$ 平面 ${PBC}$ . 试判断四面体 ${EBCD}$ 是否为鳖臑，若是,写出其每个面的直角 (只需写出结论); 若不是,请说明理由
        \item 若 ${PD} = {CD} = 2\sqrt{2}$ ， ${BC} = 2$ ，点 $F$ 在 ${BD}$ 上运动，试求 $\bigtriangleup  {EFC}$ 面积的最小值.
    \end{enumerate}

    \begin{jieda}
        \begin{enumerate}
            \item $\because {PD} \bot$ 面 ${ABCD},{AD} \subset$ 面 ${ABCD},\therefore {PD} \bot  {AD}$ .又因为 ${ABCD}$ 为矩形, $\therefore {AD} \bot  {DC}$，且${DC},{PD} \subset$ 面 ${PCD},{DC} \cap  {PD} = D$. 故${AD} \bot$ 面 ${PDC},{DE} \subset$ 平面 ${PDC}$，因此${AD} \bot  {DE}$ ,
            
            又 $\because {DE}$ 为 $Rt{ \bigtriangleup  }{PCD}$ 斜边中线，且 ${PD} = {CD}$ ， $\therefore {DE} \bot  {PC}$
            
            $\because {BC}//{AD},\therefore {BC} \bot  {DE}$ .
            
            又因为$BC,{PC} \subset$ 面 ${PBC}$，${BC} \cap  {PC} = C$，所以${DE} \bot$ 面 ${PBC}$ .
            
            四面体 ${EBCD}$ 是鳖臑，$\angle {CED},\angle {BED},\angle {BCE},\angle {DCB}$ 为直角.

            \item 过 $F$ 作 ${PC}$ 的垂线交 ${PC}$ 于 $H$ 点,分析可知当 ${FH}$ 为 ${BD}$ 和 ${PC}$ 公垂线段时, ${FH}$ 最小,设公垂线段长为 $d$ ,则 $\bigtriangleup {EFC}$ 面积最小
            
            \begin{center}
            \includegraphics[width=0.3\textwidth]{bodyxb/src/bo_d5abfif7aajc738167kg_3_328_1391_409_388_0.jpg}
            \end{center}
            
            如图建立所示的空间直角坐标系,则 $B\left( {2,2\sqrt{2},0}\right) ,D\left( {0,0,0}\right)$ ,
            
            ${DB} = \left( {2,2\sqrt{2},0}\right) ,\vv{PC} = \left( {0,2\sqrt{2}, - 2\sqrt{2}}\right)$
            
            则法向量$\vv{n} = \left( {-\sqrt{2},1,1}\right), E\left( {0,\sqrt{2},\sqrt{2}}\right) ,\vv{EB} = \left( {2,\sqrt{2}, - \sqrt{2}}\right)$ ,
            
            则 $d = \dfrac{\left| \vv{EB} \cdot  \vv{n}\right| }{\left| \vv{n}\right| } = \dfrac{2\sqrt{2}}{2} = \sqrt{2},{CE} = \dfrac{1}{2}{PC} = 2$
            
            则 $\bigtriangleup {EFC}$ 面积的最小值为 $\dfrac{1}{2} \cdot  {EC} \cdot  \sqrt{2} = \dfrac{1}{2} \times  2 \times  \sqrt{2} = \sqrt{2}$ .

        \end{enumerate}
    \end{jieda}
    
\end{enumerate}

\clearpage

\section*{锥体——棱锥和圆锥}
A组：

\begin{enumerate}
    \item 圆锥的母线和旋转轴的位置关系是\blkda{相交}.
    
    \item 已知 $M =$ \{正多面体\}， $N =$ \{多面体\}， $R =$ \{凸多面体\}， $Q =$ \{棱长相等的三棱锥\}， 则集合 $M\text{ 、 }N\text{ 、 }R\text{ 、 }Q$ 之间的关系是\blkda[3]{$Q \subset M \subset R \subset N$}.

    \item 底面是正三角形，且每个侧面是等腰三角形的三棱锥是 \dotfill{(\kh{D})}  
    \xx{一定是正三棱锥}{一定是正四面体}{不是斜三棱锥}{可能是斜三棱锥}
    \begin{jieda}
        5个棱长为$a$，一个棱长为$b$的三棱锥
    \end{jieda}
    
    \item 下列命题中假命题的是 \dotfill{(\kh{D})}  
    \xx{多面体的面数最少是 4 个}{正多面体有且只有五种}{四面体都是三棱锥}{五面体就是三棱柱}
    
    \item 若三棱锥的顶点在底面的射影恰是底面三角形的内心，则三棱锥的\dotfill{(\kh{D})}
    
    \xx{三条侧棱长一定相等}{三条侧棱两两垂直}{三条侧棱与底面所成的角相等}{三个侧面与底面所成的二面角相等}
\end{enumerate}

B组：

\begin{enumerate}
    \item “棱锥侧棱长相等”是 “棱锥底面多边形有一个外接圆” 的 \dotfill{(\kh{B})}
    \xx{充要条件}{充分非必要条件}{必要非充分条件}{既非充分也非必要条件}
    \begin{jieda}
        棱锥侧棱长相等$\Longleftrightarrow$顶点在底面投影点到底面多边形的顶点距离相等$\Longrightarrow$棱锥底面多边形有一个外接圆
    \end{jieda}
    
    \item 若四棱锥的四个侧面与底面所成的角都相等，则其底面四边形必是\dotfill{(\kh{C})}
    \xx{矩形}{菱形}{圆外切四边形}{圆内接四边形}
    \begin{jieda}
        四棱锥的四个侧面与底面所成的角都相等$\Longleftrightarrow$顶点在底面的射影到地面的边的距离相等
    \end{jieda}

    \item 若正棱锥的底面边长和侧棱长相等，则该棱锥一定不是 \dotfill{(\kh{D})}  
    \xx{三棱锥}{四棱锥}{五棱锥}{六棱锥}
    
    \item 在侧面为正三角形的正棱锥, 其侧面与底面所成的二面角的大小是\blkda[8]{$\arccos \left(\dfrac{1}{3}\right)$或$\arccos \left(\dfrac{\sqrt{3}}{3}\right)$或$\arctan \left( {3 - \sqrt{5}}\right)$}.
    \begin{jieda}
        设底面边长为$a$，侧面和底面的二面角大小为$\theta$（锐角），根据$cos\theta = \dfrac{S}{S'}$，$S' = \dfrac{\sqrt{3}}{4}a^2$
        \begin{dingautolist}{172}
            \item 正三棱锥：$S = \dfrac{\sqrt{3}}{4}a^2 \times \dfrac{1}{3}$，故$\theta = arccos\left(\dfrac{1}{3}\right)$
            \item 正四棱锥：$S = a^2 \times \dfrac{1}{4}$，故$\theta = arccos\left(\dfrac{\sqrt{3}}{3}\right)$
            \item 正五棱锥：$S = \dfrac{a^2}{4} \times cot\ang{36} = \dfrac{a^2}{4} \times \sqrt{1+\dfrac{2\sqrt{5}}{5}}$，故$\theta = arccos\left(\sqrt{\dfrac{1}{3}+\dfrac{2\sqrt{5}}{15}}\right) = \arctan \left( {3 - \sqrt{5}}\right)$
        \end{dingautolist}
    \end{jieda}
    
    \item 若正三棱锥 $S-ABC$ 的侧棱与底面边长相等，$E$、$F$ 分别是 $SC$、$AB$ 的中点， 则直线 $EF$ 与 $SA$ 所成的角为\blkda{$\dfrac{\pi}{4}$}.

    \item 下列说法正确的是\dotfill{(\kh{C})}
    \xx{有两个面平行，其他各个面都是平行四边形的多面体是棱柱}{过圆锥顶点的所有截面中，轴截面面积最大}{有一个面是平行四边形的棱锥一定是四棱锥}{有两个面互相平行且相似, 其他各个面都是梯形的多面体是棱台}
    \begin{jieda}
        对于 ${A}$ ,如图 1 所示,上下底面平行,各个面都是平行四边形，此几何体不是棱柱,故 ${A}$ 错误
        
        对于 ${B}$ ,过圆锥顶点的截面为等腰三角形,且两腰长为母线长 $l$ ,
        设该等腰三角形顶角为 $\theta$ ,则截面三角形面积为 $S = \dfrac{1}{2}{l}^{2}\sin \theta$，显然当 $\theta  = \dfrac{\pi }{2}$ ,面积 ${S}$ 最大，故当圆锥的轴截面三角形顶角大于 $\dfrac{\pi }{2}$ 时,圆锥的轴截面面积不一定是最大的,故 B 错误
        
        对于 ${C}$ ,棱锥侧面全为三角形,有一个面是平行四边形,则此面为底面,
        所以该棱锥为四棱锥,故 ${C}$ 正确
        
        对于 ${D}$ ,棱台是棱锥截得的,侧棱的延长线要交于同一点.
        有两个面平行且相似, 其他各个面都是梯形的多面体,
        不能保证侧棱的延长线交于同一点,因此该多面体不一定是棱台,故 ${D}$ 错误.
        
        反例 1: 如图 2 所示的几何体中,四边形 ${ABCD}$ 和四边形 ${A}_{1}{B}_{1}{C}_{1}{D}_{1}$ 均为矩形,且 ${AB} = 4,{BC} = 2$, ${A}_{1}{B}_{1} = 1,{B}_{1}{C}_{1} = 2$ ,且平面 ${ABCD}//$ 平面 ${A}_{1}{B}_{1}{C}_{1}{D}_{1}$ ,且 ${AB}//{A}_{1}{B}_{1},{CD}//{D}_{1}{C}_{1}$ ,则该几何体满足 ${D}$ ,但不是棱台;
        
        反例 2: 如图 3 所示,由两个棱台组合而成的几何体满足 ${D}$ ,但这个几何体不是棱台.
        \begin{center}
            \includegraphics[width=0.5\linewidth]{bodyxb//src/zhuiti1.png}
        \end{center}
    \end{jieda}

    \item 正三棱锥 $S-ABC$ 的高 ${{SO}} = h$ ，斜高 ${{SM}} = l$ ，点 ${P}$ 在 ${{SO}}$ 上且分 ${{SO}}$ 所成的比是 $1 : 2$ ，则过 ${P}$ 点且平行于底面的截面面积是\blkda{$\dfrac{\sqrt{3}(l^2-h^2)}{3}$}.
    \begin{jieda}
        $\dfrac{S'}{S} = \left(\dfrac{h'}{h}\right)^2 = \dfrac{1}{9}$

        $S = (2\sqrt{3}\times\sqrt{l^2-h^2})^2 \times \dfrac{\sqrt{3}}{4} = 3\sqrt{3}(l^2-h^2)$，故$S' = \dfrac{\sqrt{3}(l^2-h^2)}{3}$
    \end{jieda}

    \item 若矩形 ${{ABCD}}$ 的边长分别为 $a,b\left( {\dfrac{b}{2} < a < {2b}}\right) ,{E}$ 是 ${{DC}}$ 的中点，把矩形沿 ${{AE}}$ 、 BE 折成一个三棱锥的三个侧面(C、D 重合)，则侧面 ABC 与底面 ABE 所成的二面角的正弦值是\blkda{$\dfrac{a}{2b}$或$ \dfrac{b}{2a}$}.
    \begin{jieda}
        因为${EC} \bot  {BC},{ED} \bot  {AD},C,D$ 重合，所以${EC} \bot$ 平面 ${ABC}$
        
        取 ${AB}$ 中点 $H$ ,连结 ${CH}$ , ${EH}$ ,则 ${EH} \bot  {AB}$ , ${CH} \bot  {AB}$，故$\angle {CHE}$ 就是所成二面角的平面角
        
        $\sin \angle {CHE} = \dfrac{EC}{EH} = \dfrac{\dfrac{a}{2}}{b} = \dfrac{a}{2b}$或$\sin \angle {CHE} = \dfrac{EC}{EH} = \dfrac{\dfrac{b}{2}}{a} = \dfrac{b}{2a}$

        \begin{center}
            \includegraphics[width=0.25\linewidth]{bodyxb//src/zhuiti2.png}
        \end{center}
    \end{jieda}

    \item 如图，某圆锥的轴截面 ${ABC}$ 是等边三角形，点 $D$ 是线段 ${AB}$ 的中点，点 $E$ 在底面圆的圆周上,且 $\overset{\frown}{BE}$ 的长度等于 $\overset{\frown}{CE}$ 的长度,则异面直线 ${DE}$ 与 ${BC}$ 所成角的余弦值是\blkda{$\dfrac{\sqrt{2}}{4}$}.
    
    \begin{center}
    \includegraphics[width=0.2\textwidth]{bodyxb/src/bo_d5ar5ef7aajc73819veg_1_252_764_282_278_0.jpg}
    \end{center}
    \begin{jieda}
        过点 ${A作AO}\bot {BC}$ 于点 $O$ ,过点 ${A作DG}\bot {BC}$ 于点 $G$ ,取 ${AO}$ 的中点 $F$ ,连接 ${GE}$ 、 ${OE}$ 、 ${EF}$，则 ${DF}//{BC}$ ,且 ${DF} = \dfrac{1}{2}{BC}$ ,所以 $\angle {DEF}$ (或其补角)就是异面直线 ${DE}$ 与 ${BC}$ 所成的角
        
        设圆锥的底面半径为 2,则 ${DF} = 1,{OE} = 2,{AO} = 2\sqrt{3}$ ,所以 ${DG} = {OF} = \sqrt{3}$ ,
        在$Rt\triangle {GOE}$ 中, ${GO} = 1,{OE} = 2$ ,所以 ${GE} = \sqrt{G{O}^{2} + O{E}^{2}} = \sqrt{5}$ 
        
        在$Rt{\triangle GDE}$ 中, ${GE} = \sqrt{5},{DG} = \sqrt{3}$ ,所以 ${DE} = \sqrt{G{D}^{2} + G{E}^{2}} = 2\sqrt{2}$，在$Rt {\triangle FOE}$ 中, ${FO} = \sqrt{3},{OE} = 2,{FE} = \sqrt{F{O}^{2} + O{E}^{2}} = \sqrt{7}$
        
        所以在 ${\triangle DFE}$ 中,满足 $D{F}^{2} + F{E}^{2} = D{E}^{2}$ ,所以 $\angle {DFE} = {90}^{ \circ  }$，所以 $\cos \angle {DEF} = \dfrac{DF}{DE} = \dfrac{1}{2\sqrt{2}} = \dfrac{\sqrt{2}}{4}$
    \end{jieda}


    \item 如图，某几何体平面展开图由一个等边三角形和三个等腰直角三角形组合而成， $E$ 为 ${BC}$ 的中点，则在原几何体中，异面直线 ${AE}$ 与 ${CD}$ 所成角的余弦值为\blkda{$\dfrac{\sqrt{6}}{6}$}.
    
    \begin{center}
    \includegraphics[width=0.2\textwidth]{bodyxb/src/bo_d5ar5ef7aajc73819veg_1_1139_1166_301_307_0.jpg}
    \end{center}
    \begin{jieda}
        因几何体平面展开图由一个等边三角形和三个等腰直角三角形组合而成,于是得原几何体是正三棱锥 $D - {ABC}$,其中 ${DA},{DB},{DC}$ 两两垂直,且 ${DA} = {DB} = {DC}$ ,取 ${BD}$ 中点 $F$ ,连接 ${EF}$ , ${AF}$ ,如图所示

        \begin{center}
            \includegraphics[width=0.25\linewidth]{bodyxb//src/zhuiti3.png}
        \end{center}
        
        因 $E$ 为 ${BC}$ 的中点,则有 ${EF}//{CD}$ ,因此, $\angle {AEF}$ 是异面直线 ${AE}$ 与 ${CD}$ 所成角或其补角

        
        令 ${DB} = 2$ ,则 ${EF} = \dfrac{1}{2}{DC} = 1$ , $Rt \bigtriangleup  {ADF}$ 中, ${AF} = \sqrt{A{D}^{2} + D{F}^{2}} = \sqrt{5}$ ,
        正 $\bigtriangleup {ABC}$ 中, ${AE} = \dfrac{\sqrt{3}}{2}{AB} = \sqrt{6}$ ,于是有： $A{F}^{2} + E{F}^{2} = 6 = A{E}^{2}$ ,即 $\angle {AFE} = {90}^{ \circ  }$ , $\cos \angle {AEF} = \dfrac{EF}{AE} = \dfrac{\sqrt{6}}{6}$ 
        
        所以异面直线 ${AE}$ 与 ${CD}$ 所成角的余弦值为 $\dfrac{\sqrt{6}}{6}$ .
    \end{jieda}

    \item 如图所示的几何体是由一个圆柱挖去一个以圆柱上底面为底面，下底面圆心为顶点的圆锥而得到的几何体，现用一个竖直的平面去截这个几何体，则截面图形可能是\dotfill{(\kh{D})}
    \xx{(2) (5)}{(1) (3)}{(2) (4)}{(1) (5)}
        
    \begin{center}
    \includegraphics[width=0.6\textwidth]{bodyxb/src/bo_d5ar5ef7aajc73819veg_1_241_1429_740_188_0.jpg}
    \end{center}
    
    \item 若棱长为 $a$ 的正四面体中，高为 ${H}$ ，斜高为 $h$ ，相对棱间的距离为 $d$ ，则 $a.{H}.h.d$ 的大小关系是\dotfill{(\kh{B})}  \xx{$a > H > h > d$}{$a > h > H > d$}{$a > h > d > H$}{$a > d > h > H$}
    \begin{jieda}
        $h = \dfrac{\sqrt{3}}{2}a, H = \dfrac{\sqrt{6}}{3}a, d = \dfrac{\sqrt{2}}{2}a$
    \end{jieda}

    \item 如果正三棱锥的高为 $\sqrt{3}$ ，侧棱长为 $\sqrt{7}$ ，那么侧面与底面所成的二面角\dotfill{(\kh{B})}  
    \xx{大小为 ${30}^{ \circ  }$}{大小为 ${60}^{ \circ  }$}{正弦为 $\dfrac{\sqrt{21}}{7}$}{余弦为 $\dfrac{\sqrt{21}}{7}$}

    \item 如图，已知底面是正方形的四棱锥，其一条侧棱与底面垂直，它的长与底面边长相等，长度均为 1 ，那么该棱锥中最长的棱长是 \dotfill{(\kh{B})}
    
    \begin{center}
    \includegraphics[width=0.2\textwidth]{bodyxb/src/bo_d5ar5ef7aajc73819veg_2_1089_199_288_237_0.jpg}
    \end{center}
    
    \xx{$\sqrt{2}$}{$\sqrt{3}$}{$2\sqrt{2}$}{3}
    
    \item 如果正三棱锥的侧面都是直角三角形，则侧棱与底面所成的角 $\theta$ 应属于\dotfill{(\kh{C})}  
    
    \xx{$\left( {0,\dfrac{\pi }{6}}\right)$}{$\left( {\dfrac{\pi }{4},\dfrac{\pi }{3}}\right)$}{$\left( {\dfrac{\pi }{6},\dfrac{\pi }{4}}\right)$}{$\left( {\dfrac{\pi }{3},\dfrac{\pi }{2}}\right)$}
    \begin{jieda}
        该三棱锥可沿着立方体的三等分平面切出，故可知$\theta = arcsin\left(\dfrac{\sqrt{3}}{3}\right)$
    \end{jieda}

    \item 已知正四面体 ${ABCD}$ 中，E、F 分别为 ${BC}$ 、 ${AD}$ 的中点，连 ${AE}$ 、 ${CF}$
    
    \begin{center}
    \includegraphics[width=0.2\textwidth]{bodyxb/src/bo_d5ar5ef7aajc73819veg_2_999_733_332_281_0.jpg}
    \end{center}

    \begin{enumerate}
        \item 求异面直线 ${AE}$ 与 ${CF}$ 所成角的大小
        \item 求$CF$与平面 ${BCD}$ 所成角的大小
    \end{enumerate}
    \begin{jieda}
        构建立方体$AMDN-PBCQ$，使得$A,B,C,D$为立方体的顶点，易知$AM, AN, AP$两两垂直，故以$A$为原点，$\vv{AM}, \vv{AN}, \vv{AP}$为$x, y, z$轴正方向建立空间直角坐标系

        不妨设棱长为2，故$B(2,0,2), C(0,2,2),D(2,2,0)$，因此$E(1,1,2), F(1,1,0)$
        \begin{enumerate}
            \item $\vv{AE} = (1,1,2)$, $\vv{CF} = (1,-1,-2)$，故$cos<\vv{AE}, \vv{CF}> = \dfrac{-4}{6}$，因此异面直线 ${AE}$ 与 ${CF}$ 所成角的大小为$arccos\left(\dfrac{2}{3}\right)$
            \item $\vv{CF} = (1,-1,-2)$，易知$\vv{n} = (1,1,1)$，则$cos<\vv{CF}, \vv{n}> = \dfrac{-2}{3\sqrt{2}} = -\dfrac{\sqrt{2}}{3}$

            故$CF$与平面 ${BCD}$ 所成角的大小为$arcsin\left(\dfrac{\sqrt{2}}{3}\right)$
        \end{enumerate}
    \end{jieda}


    \item 在梯形 ${ABCD}$ 中， ${AD}//{BC},\angle {ABC} = \dfrac{\pi }{2},{AB} = a,{AD} = {3a}$ 且 $\angle {ADC} =  \arcsin \dfrac{\sqrt{5}}{5}$ ,又 ${{PA}} \bot$ 平面 ${{ABCD}},{{PA}} = a$ . 求:
    \begin{enumerate}
        \item 二面角 $P - CD - A$ 的大小
        \item 点 $A$ 到平面 ${PBC}$ 的距离
    \end{enumerate}
    
    \begin{center}
    \includegraphics[width=0.3\textwidth]{bodyxb/src/bo_d5ate7ref24c73bimur0_0_256_661_356_315_0.jpg}
    \end{center}
    \begin{jieda}
        如图所示:
        
        \begin{center}
        \includegraphics[width=0.3\textwidth]{bodyxb/src/bo_d5ate7ref24c73bimur0_0_273_1710_406_354_0.jpg}
        \end{center}
        
        过 $A$ 作 ${AE} \bot  {CD}$ ,连接 ${PE}$，因为 ${PA} \bot$ 平面 ${ABCD},{CD} \subset$ 平面 ${ABCD}$，所以 ${PA} \bot  {CD}$ ,又 ${PA} \cap  {AE} = A$， 所以 ${CD} \bot$ 平面 ${PAE}$
        
        所以 $\angle {PEA}$ 二面角 $P - {CD} - A$ 的平面角
        
        因为 $\angle {ADC} = \arcsin \dfrac{\sqrt{5}}{5}$ ，所以 ${AE} = {AD} \times  \sin \angle {ADC} = \dfrac{3\sqrt{5}}{5}a$ ，所以 $\tan \angle {PEA} = \dfrac{AP}{AE} = \dfrac{\sqrt{5}}{3}$，所以 $\angle {PEA} = \arctan \dfrac{\sqrt{5}}{3}$，即二面角 $P - {CD} - A$ 的大小 $\arctan \dfrac{\sqrt{5}}{3}$ .
        
        (2)如图所示:
        
        \begin{center}
        \includegraphics[width=0.3\textwidth]{bodyxb/src/bo_d5ate7ref24c73bimur0_1_291_1012_364_317_0.jpg}
        \end{center}
        
        过 $A$ 作 ${AF} \bot  {PB}$，因为 $\angle {ABC} = {90}^{ \circ  }$ ， 所以 ${AB} \bot  {CB}$
        
        因为 ${PA} \bot$ 平面 ${ABCD},{BC} \subset$ 平面 ${ABCD}$，所以 ${PA} \bot  {BC}$ ,又 ${PA} \cap  {AB} = A$，所以 ${BC} \bot$ 平面 ${PAE}$ ,又 ${BC} \subset$ 平面 ${PBC}$，所以 ${PAB} \bot$ 平面 ${PBC}$
        
        又平面 ${PAB} \cap$ 平面 ${PBC} = {PB}$，所以 ${AF} \bot$ 平面 ${PBC}$，所以 ${AF}$ 为点 $A$ 到平面 ${PBC}$ 的距离
        
        在 $R{t}{\bigtriangleup PAB}$ 中, ${AF} = \dfrac{{PA} \times  {AB}}{PB} = \dfrac{a \times  a}{\sqrt{2}a} = \dfrac{\sqrt{2}}{2}a$ ，所以点 $A$ 到平面 ${PBC}$ 的距离为 $\dfrac{\sqrt{2}}{2}a$ .
    \end{jieda}

\end{enumerate}

C组：

\begin{enumerate}
    \item 如图，水平的广场上有一盏路灯挂在高 $9{m}$ 的电线杆顶上，记电线杆的底部为点 ${A}$ . 把路灯看作一个点光源，身高 ${1.5}{m}$ 的女孩站在离点 ${A}$ $5{m}$ 的点 $B$ 处，若女孩沿 ${AB}$ 方向前行 $5{m}$ 到达点 $D$ ，此时 $B$ 为 ${AD}$ 的中点，然后从点 $D$ 出发沿着以 ${BD}$ 为对角线的正方形走一圈，则女孩走一圈时头顶影子的轨迹围成图形的面积为\blkda{18} ${{m}}^{2}$ .

    \begin{center}
    \includegraphics[width=0.2\textwidth]{bodyxb/src/bo_d5ar5ef7aajc73819veg_3_238_586_328_310_0.jpg}
    \end{center}

    \begin{jieda}
        设头顶轨迹面积为$S'$，影子轨迹为$S$，则根据椎体的性质知$\dfrac{S'}{S} = \left(\dfrac{h'}{h}\right)^2 = \left(\dfrac{9-1.5}{9}\right)^2 = \dfrac{25}{36}$

        $S' = \dfrac{25}{2}m^2$，故$S = 18m^2$
    \end{jieda}
    
    \item 给定四面体 ${ABCD}$ . 平面 $a$ 满足: \ding{172} ${A}$ 、 $B$ 、 $C$ 、 $D$ 四个点均不在平面 $a$ 上， 也不在 $a$ 的同侧;  \ding{173} 若平面 $a$ 与四面体 ${ABCD}$ 的棱有公共点，则该公共点一定是此棱的中点或两个三等分点之一. 设 ${A}\text{ 、 }B\text{ 、 }C\text{ 、 }D$ 四个点到平面 $a$ 的距离分别为 ${d}_{i}\left( {i = 1,2,3,4}\right)$ ，那么 ${d}_{i}$ 的所有不同值的个数组成的集合为\dotfill{(\kh{B})}  
    \xx{$\{ 1,2,3,4\}$}{$\{ 1,2,3\}$}{$\{ 1,2\}$}{$\{ 1\}$}
\end{enumerate}

\section*{锥体的体积和表面积}

A 组：
\begin{enumerate}
    \item 给出下列命题:
    \begin{dingautolist}{172}
        \item 底面是正多边形的棱锥是正棱锥
        \item 侧棱都相等的棱锥是正棱锥
        \item 侧棱和底面成等角的棱锥是正棱锥
        \item 侧面和底面所成二面角都相等的棱锥是正棱锥
        \item 侧棱和底面成等角、侧面和底面所成二面角都相等的棱锥是正棱锥
    \end{dingautolist}
    其中正确的命题是\blkda{\ding{176}}.
    
    \item 若正四面体的棱长为 $a$ ，则相邻两个面的夹角的余弦是\blkda{$\dfrac{1}{3}$}.
    
    \item 棱长都是 3 的三棱锥的高等于\blkda{$\sqrt{6}$}.

    \item 边长为 $3\sqrt{2}$ 的正四面体的一个顶点到对应底面的距离为\blkda{$2\sqrt{3}$}.
    
    \item 已知圆锥的底面半径为 $2\sqrt{2}$ ，其侧面展开图为一个半圆，则该圆锥的母线长为\blkda{$4\sqrt{2}$}.
    \begin{jieda}
        设圆锥的母线长为 $l$ ,由于圆锥底面圆的周长等于扇形的弧长,则 ${\pi l} = {2\pi } \times  2\sqrt{2}$，解得 $l = 4\sqrt{2}$
    \end{jieda}
\end{enumerate}

B 组：

\begin{enumerate}
    \item 若一个圆锥的侧面展开图是圆心角为 ${120}^{ \circ  }$ ，半径为 1 的扇形，则这个圆锥的表面积与侧面积的比是\blkda{$4:3$}.
    \begin{jieda}
        设圆锥底面圆的半径为 $r$ ,则底面圆的周长为 ${2\pi r}$ ,即展开后的扇形弧长为 ${2\pi r}$ 
        
        又扇形的圆心角为$\ang{120}$,半径为 1,
        所以 $\dfrac{{120}^{ \circ  }}{{360}^{ \circ  }} \times  {2\pi } \times  1 = {2\pi r}$ ,所以 $r = \dfrac{1}{3}$
        
        故圆锥的侧面积为 $\dfrac{1}{2} \times  1 \times  \dfrac{2\pi }{3} = \dfrac{\pi }{3}$，表面积为 $\dfrac{\pi }{3} + \pi  \times  {\left( \dfrac{1}{3}\right) }^{2} = \dfrac{4\pi }{9}$ 
        
        所以这个圆锥的表面积与侧面积的比为 $\dfrac{\dfrac{4\pi }{9}}{\dfrac{\pi }{3}} = \dfrac{4}{3}$ ,
        即这个圆锥的表面积与侧面积的比是 $4 : 3$ .
    \end{jieda}

    \item \begin{dingautolist}{172}
        \item 三棱锥的四个面都可能为直角三角形
        \item 正四面体在某平面的射影可能为正方形
        \item 三棱锥 $A - {BCD}$ 中， $A$ 在 $\bigtriangleup  {BCD}$ 的射影为 $\bigtriangleup  {BCD}$ 的重心，则 $B$ 在 $\bigtriangleup {ACD}$ 的射影也是 $\bigtriangleup {ACD}$ 的重心
        \item 四棱锥 $P - {ABCD}$ 中若各侧棱与底面所成角均相等，则四边形 ${ABCD}$ 为菱形
    \end{dingautolist}
    其中正确命题是\blkda{\ding{172}\ding{173}}.
    \begin{jieda}
        \begin{dingautolist}{172}
            \item 鳖臑
            \item 正四面体可以补成正方体
            \item 正三棱锥$A - {BCD}$满足$A$ 在 $\bigtriangleup  {BCD}$ 的射影为 $\bigtriangleup  {BCD}$ 的重心，随着棱锥的高变化，$B$ 在 $\bigtriangleup {ACD}$ 的射影不一定是 $\bigtriangleup {ACD}$ 的重心
            \item 设$P$在底面射影为$O$，则$\triangle POA,\triangle POB,\triangle POC,\triangle POD$全等，即$OA = OB=OC=OD$，四边形 ${ABCD}$ 不一定为菱形
        \end{dingautolist}
    \end{jieda}
    
    \item 已知正三棱锥 $P - {ABC}$ 的底面边长为 6，点 $P$ 到底面 ${ABC}$ 的距离为 3，则三棱锥的表面积是\blkda{$27\sqrt{3}$}.
    \begin{jieda}
        由题意可知底面三角形的中心到底面三角形的边的距离为 $\dfrac{1}{3} \times  \dfrac{\sqrt{3}}{2} \times  6 = \sqrt{3}$ ,
        所以正三棱锥的斜高为 $\sqrt{{3}^{2} + {\left( \sqrt{3}\right) }^{2}} = 2\sqrt{3}$ ,
        所以这个正三棱锥的侧面积为 $3 \times  \dfrac{1}{2} \times  6 \times  2\sqrt{3} = {18}\sqrt{3}$ ,底面积为 $\dfrac{1}{2} \times  {6}^{2} \times  \sin {60}^{ \circ  } = 9\sqrt{3}$ ,
        所以正三棱锥 $P - {ABC}$ 的表面积为 ${18}\sqrt{3} + 9\sqrt{3} = {27}\sqrt{3}$ .
    \end{jieda}

    \item 直角坐标系 ${xOy}$ 内有点 $P\left( {-2,0}\right) ,Q\left( {0, - 2}\right)$ ，将 $\bigtriangleup  {POQ}$ 绕 $x$ 轴旋转一周，则所得几何体的体积为\blkda{$\dfrac{8\pi}{3}$}.
    \begin{jieda}
        由题意可知 $\left| {OP}\right|  = \left| {OQ}\right|  = 2$ ,将 $\bigtriangleup {POQ}$ 绕 $x$ 轴旋转一周,则所得的几何体是以 2为底面半径, 2 为高的圆锥,所以所得几何体的体积为 $\dfrac{1}{3}\pi  \times  {2}^{2} \times  2 = \dfrac{8\pi }{3}$ ,
    \end{jieda}

    \item 已知四棱锥$P-ABCD$, $PA \bot $平面$ABCD$, 且 $ABCD$ 是矩形，$PB$、$PC$、$PD$ 的长分别为 $\sqrt{5},\sqrt{17},\sqrt{13}$ ,则 ${P}$ 到 ${{BD}}$ 的距离是\blkda{$2$}.
    \begin{jieda}
        $\left\{\begin{array}{l}
             PA^2 + AB^2 = 5  \\
             PA^2 + AC^2 = 17  \\
             PA^2 + AD^2 = 13  \\
             AB^2+AD^2 = AC^2
        \end{array}\right. \Longrightarrow \left\{\begin{array}{l}
             PA = 1  \\
             AB = 2  \\
             AD = 2\sqrt{3}
        \end{array}\right.$，$d = \sqrt{1+3} = 2$
    \end{jieda}
    
    \item 如图，在正四棱锥 $P - {ABCD}$ 中， ${PA} = 4$ ， ${AB} = 2\sqrt{2}$ . 从 $A$ 拉一条细绳绕过侧棱 ${PB},{PC},{PD}$ 回到 $A$ 点，则细绳的最短长度为\blkda{$3\sqrt{7}$}.
    \begin{center}
    \includegraphics[width=0.2\textwidth]{bodyxb/src/bo_d5b4icbef24c73bio840_1_238_208_302_324_0.jpg}
    \end{center}

    \begin{jieda}
        如图,将侧面 ${PAB}$ ,侧面 ${PBC}$ ,侧面 ${PCD}$ ,侧面 ${PDA}$ 展开到一个平面内.

        \begin{center}
        \includegraphics[width=0.25\linewidth]{bodyxb//src/zhuiti4.png}
        \end{center}
        
        由题意可知 ${PA} = {PB} = {PC} = {PD} = P{A}^{\prime } = 4,{AB} = {BC} = {CD} = D{A}^{\prime } = 2\sqrt{2}$ ,
        设 $\angle {APB} = \angle {BPC} = \angle {CPD} = \angle {DP}{A}^{\prime } = \theta$ ,则 $\cos \theta  = \dfrac{{16} + {16} - 8}{2 \times  4 \times  4} = \dfrac{3}{4}$ ,
        所以 $\cos {2\theta } = 2{\cos }^{2}\theta  - 1 = \dfrac{1}{8}$ ,所以 $\cos {4\theta } = 2{\cos }^{2}{2\theta } - 1 =  - \dfrac{31}{32}$ .
        由余弦定理可得 ${A}^{\prime }{A}^{2} = {4}^{2} + {4}^{2} - 2 \times  4 \times  4 \times  \left( {-\dfrac{31}{32}}\right)  = {63}$ ,
        则 ${A}^{\prime }A = 3\sqrt{7}$ ,即细绳的最短长度为 $3\sqrt{7}$ .
    \end{jieda}

    

    \item 在一张硬纸上，挖去一个半径为 $\sqrt{3}$ 的圆洞，然后把此洞套在一个底面边长为4，高为 6 的正三棱锥上，并使纸面与棱锥底面平行，则当正三棱锥恰好卡在圆洞内时，穿过这张纸面的部分棱锥的高等于\blkda{$\dfrac{9}{2}$}.
    \begin{jieda}
        如图所示,假设硬纸上的圆洞刚好卡在 ${B}^{\prime }{C}^{\prime }{D}^{\prime }$处. 
        \begin{center}
            \includegraphics[width=0.25\linewidth]{bodyxb//src/zhuiti7.png}
        \end{center}
            
        设正三棱锥 $A - {BCD}$ 的顶点 $A$ 在平面 ${BCD}$ 上的射影为 ${A}^{\prime }$ ,在平面 ${B}^{\prime }{C}^{\prime }{D}^{\prime }$ 上的射影为 $O$ . 连接 $B{A}^{\prime }\text{、}{B}^{\prime }O$ 并延长分别交 ${CD}\text{、}{C}^{\prime }{D}^{\prime }$ 于 $E\text{、}{E}^{\prime }$ 点,那么平面 ${B}^{\prime }{C}^{\prime }{D}^{\prime }//$ 平面${BCD}$ ,所以 $\dfrac{{B}^{\prime }{E}^{\prime }}{BE} = \dfrac{{B}^{\prime }{C}^{\prime }}{BC},{B}^{\prime }{E}^{\prime } = \dfrac{3}{2}{B}^{\prime }O,{BE} = \dfrac{3}{2}B{A}^{\prime }$ ,即$\dfrac{{B}^{\prime }O}{B{A}^{\prime }} = \dfrac{{B}^{\prime }{C}^{\prime }}{BC}$ . 又因为 $B{A}^{\prime } = \dfrac{4\sqrt{3}}{3},{B}^{\prime }O = \sqrt{3},{BC} = 4$ ,所以${B}^{\prime }{C}^{\prime } = 3$ ,又 $\dfrac{AO}{A{A}^{\prime }} = \dfrac{{B}^{\prime }{C}^{\prime }}{BC}$ ,所以 ${AO} = \dfrac{{B}^{\prime }{C}^{\prime }}{BC} \cdot  A{A}^{\prime } = \dfrac{9}{2}$ .
    \end{jieda}

    \item 在四棱锥 $P - {ABCD}$ 中，底面 ${ABCD}$ 是边长为 2 的正方形， ${PA} \bot$ 平面 ${ABCD}$ ， 且 ${PA} = 2$ . 若点 $E,F,G$ 分别为棱 ${AB},{AD},{PC}$ 的中点，则正确的是\blkda{\ding{172}\ding{173}\ding{175}}.
    \begin{dingautolist}{172}
        \item ${AG} \bot$ 平面 ${PBD}$ 
        \item 直线 ${FG}$ 和直线 ${AB}$ 所成的角为 $\dfrac{\pi }{4}$
        \item 当点 $T$ 在平面 ${PBD}$ 内,且 ${TA} + {TG} = 2$ 时，点 $T$ 的轨迹为一个椭圆
        \item 过点 $E,F,G$ 的平面与四棱锥 $P - {ABCD}$ 表面交线的周长为 $2\sqrt{2} + \sqrt{6}$
    \end{dingautolist}
    \begin{jieda}
    
        由题意可知因为 ${PA} \bot$ 平面 ${ABCD},{AB},{AD} \subset$ 平面 ${ABCD}$，所以 ${AP} \bot  {AB},{AP} \bot  {AD}$
        
        又底面 ${ABCD}$ 是边长为 2 的正方形,所以 ${AB} \bot  {AD}$ ,即 ${AP},{AB},{AD}$ 两两垂直
        
        以 $A$ 为原点， ${AB},{AD},{AP}$ 为 $x$ 轴， $y$ 轴， $z$ 轴建立如图所示坐标系
        
        \begin{center}
        \includegraphics[width=0.2\textwidth]{bodyxb/src/bo_d5b5t477aajc7381bj70_0_242_1352_343_334_0.jpg}
        \end{center}
        
        所以由题意 $A\left( {0,0,0}\right) ,B\left( {2,0,0}\right) ,C\left( {2,2,0}\right) ,D\left( {0,2,0}\right) ,P\left( {0,0,2}\right) ,E\left( {1,0,0}\right) ,F\left( {0,1,0}\right) ,G\left( {1,1,1}\right)$ , 所以 $\vv{AG} = \left( {1,1,1}\right) ,\vv{PB} = \left( {2,0, - 2}\right) ,\vv{PD} = \left( {0,2, - 2}\right)$
    
        \begin{dingautolist}{172}
            \item 设平面 ${PBD}$ 的法向量 $\vv{n} = \left( {x,y,z}\right)$ ,所以 $\left\{  \begin{array}{l} \vv{PB} \cdot  \vv{n} = {2x} - {2z} = 0 \\  \vv{PD} \cdot  \vv{n} = {2y} - {2z} = 0 \end{array}\right.$ ,解得 $\vv{n} = \left( {1,1,1}\right)$

            因为 $\vv{n} = 1\vv{AG}$ ,所以 ${AG} \bot$ 平面 ${PBD}$
            \item 因为 $\vv{FG} = \left( {1,0,1}\right) ,\vv{AB} = \left( {2,0,0}\right)$, 所以 $\cos  < \vv{FG},\vv{AB} >  = \dfrac{\left| \vv{FG} \cdot  \vv{AB}\right| }{\left| \vv{FG}\right| \left| \vv{AB}\right| } = \dfrac{\left| 2 \times  1\right| }{\sqrt{{1}^{2} + {1}^{2}}\sqrt{{2}^{2}}} = \dfrac{\sqrt{2}}{2}$，又因为 $< \vv{FG},\vv{AB} >  \in  \left\lbrack  {0,\dfrac{\pi }{2}}\right\rbrack$ ,所以直线 ${FG}$ 和直线 ${AB}$ 所成的角为 $\dfrac{\pi }{4}$
            \item 因为 $\vv{GB} = \left( {1, - 1, - 1}\right)$ ，所以点 $G$ 到平面 ${PBD}$ 的距离 ${d}_{1} = \dfrac{\left| \vv{n} \cdot  \vv{GB}\right| }{\left| \vv{n}\right| } = \dfrac{\left| 1 \times  1 + 1 \times  \left( -1\right)  + 1 \times  \left( -1\right) \right| }{\sqrt{{1}^{2} + {1}^{2} + {1}^{2}}} = \dfrac{\sqrt{3}}{3}$
            
            又因为 $\vv{AB} = \left( {2,0,0}\right)$ ,所以点 $A$ 到平面 ${PBD}$ 的距离 ${d}_{2} = \dfrac{\left| \vv{n} \cdot  \vv{AB}\right| }{\left| \vv{n}\right| } = \dfrac{\left| 2 \times  1\right| }{\sqrt{{1}^{2} + {1}^{2} + {1}^{2}}} = \dfrac{2\sqrt{3}}{3}$
            
            设 ${AG}$ 与平面 ${PBD}$ 交于 ${S}$ ,由 ${A}$ 得因为 ${TS} \subset$ 平面 ${PBD}$ ,所以 ${AG} \bot  {TS}$ ,
            
            所以 ${TA} + {TG} = \sqrt{T{S}^{2} + A{S}^{2}} + \sqrt{T{S}^{2} + G{S}^{2}} = \sqrt{T{S}^{2} + {\left( \dfrac{2\sqrt{3}}{3}\right) }^{2}} + \sqrt{T{S}^{2} + {\left( \dfrac{\sqrt{3}}{3}\right) }^{2}} = 2$
            
            即 ${TS}$ 为定值,所以 $T$ 的轨迹为圆
            \item 延长 ${FE}$ 与 ${CB}$ 交于点 $N$ ,延长 ${EF}$ 与 ${CD}$ 交于点 $J$ ,连接 ${NG}$ 与 ${PB}$ 交于点 $H$ ,连接 ${GJ}$ 与 ${PD}$ 交于点 $K$ ,连接 ${HE},{KF}$，则过点 $E,F,G$ 的平面与四棱锥 $P - {ABCD}$ 的截面为 ${HEFKG}$

            \begin{center}
            \includegraphics[width=0.4\textwidth]{bodyxb/src/bo_d5b5t477aajc7381bj70_1_242_423_487_330_0.jpg}
            \end{center}
            
            取 ${PB}$ 的中点为 $L$ ,则 ${LG} = \dfrac{1}{2}{BC} = {AF} = {BN}$
            
            又 ${LG}//{BN}$ ，所以 $\angle {LGH} = \angle {HNB},\angle {GLH} = \angle {NBH}$ ，所以 $\bigtriangleup  {LGH} \cong   \bigtriangleup  {BNH}$，所以 $H$ 为 ${BL}$ 中点,即 $H$ 为 ${PB}$ 靠近 $B$ 的四等分点,同理 $K$ 为 ${PD}$ 靠近 $D$ 的四等分点
            
            所以 $H\left( {\dfrac{3}{2},0,\dfrac{1}{2}}\right) ,K\left( {0,\dfrac{3}{2},\dfrac{1}{2}}\right)$，则 ${HE} = \dfrac{\sqrt{2}}{2},{EF} = \sqrt{2},{KF} = \dfrac{\sqrt{2}}{2},{GH} = {GK} = \dfrac{\sqrt{6}}{2}$ ,
            
            则截面周长为 $2\sqrt{2} + \sqrt{6}$
        \end{dingautolist}
    \end{jieda}

    \item 暑假将临，大学生小明同学准备利用假期探访名胜古迹. 已知某座山高耸入云，整体呈圆锥形，其半山腰(母线的中点) 有一座古寺，与上山入口在同一条母线上，入口和古寺通过一条盘山步道相连，且当时为了节省资金，该条盘山步道是按 “到达古寺的路程最短” 修建的. 如图, 已知该座山的底面半径 $R = 2\left( {{km}}\right)$ ，高 $h = 4\sqrt{2}\left( {{km}}\right)$ ，则盘山步道的长度为\blkda{$3\sqrt{7}km$}，其中上山(到山顶的直线距离减小)和下山(到山顶的直线距离增大) 路段的长度之比为\blkda{$5:2$}.
    
    \begin{center}
    \includegraphics[width=0.2\textwidth]{bodyxb/src/bo_d5b4icbef24c73bio840_1_232_1116_292_383_0.jpg}
    \end{center}

    \begin{jieda}
        设入口为点 ${A}$ ,古寺位置为 $B$ ,则由题意, $B$ 为母线 ${SA}$ 的中点,
        由底面圆半径为 $2{{km}}$ ,山高为 $4\sqrt{2}{{km}}$ ,则母线 ${SA} = \sqrt{{2}^{2} + {\left( 4\sqrt{2}\right) }^{2}} = 6$ ,
        底面圆周长 ${2\pi r} = {4\pi }$ ,所以展开图的圆心角 $a = \dfrac{4\pi }{6} = \dfrac{2\pi }{3}$ ,
        作圆锥侧面展开图, 如图,
        \begin{center}
            \includegraphics[width=0.25\linewidth]{bodyxb//src/zhuiti5.png}
        \end{center}
        由余弦定理可得 ${AB} = \sqrt{A{S}^{2} + S{B}^{2} - {2AS} \cdot  {SB}\cos \dfrac{2\pi }{3}} = \sqrt{{6}^{2} + {3}^{2} - 2 \times  6 \times  3 \times  \left( {-\dfrac{1}{2}}\right) } = 3\sqrt{7}$ ,
        由点 ${S}$ 向 ${AB}$ 引垂线,垂足为点 $H$ ,此时 ${SH}$ 为点 ${S}$ 和线段 ${AB}$ 上的点连线的最小值,
        即点 $H$ 为公路的最高点, ${AH}$ 段为上坡路, ${HB}$ 段为下坡路段,
        由 $\dfrac{1}{2}{SA} \cdot  {SB} \cdot  \sin \dfrac{2\pi }{3} = \dfrac{1}{2}{AB} \cdot  {SH}$ 可得, ${SH} = \dfrac{3\sqrt{21}}{7}$ ,
        所以 ${AH} = \sqrt{S{A}^{2} - S{H}^{2}} = \sqrt{{36} - \dfrac{27}{7}} = \dfrac{{15}\sqrt{7}}{7},{HB} = {AB} - {AH} = \dfrac{6\sqrt{7}}{7}$，所以上山和下山路段的长度之比为 $5 : 2$ .
    \end{jieda}
    
    \item 已知正四棱锥$P-ABCD$棱长均为 13，$M$、$N$ 分别是 $PA$、$BD$ 上的点,且 $PM$: ${MA} = {BN} : {ND} = 5 : 8$ . 
    \begin{center}
    \includegraphics[width=0.3\textwidth]{bodyxb/src/bo_d5b4icbef24c73bio840_1_186_1669_419_405_0.jpg}
    \end{center}
    
    \begin{enumerate}
        \item 求证:直线 $MN$//平面 $PBC$
        \item 求直线 MN 与底面 ABCD 所成角的正弦值.
    \end{enumerate}
    \begin{jieda}
        \begin{enumerate}
            \item 因为 $P - {ABCD}$ 是正四棱锥,所以 ${ABCD}$ 是正方形,连结 ${AN}$ 并延长交 ${BC}$ 于点 $E$ ,连结${PE}$
            
            因为 ${AD}//{BC}$ ,所以 ${EN} : {AN} = {BN} : {ND}$
            
            又因为 ${BN} : {ND} = {PM} : {MA}$ ,所以${EN} : {AN} = {PM} : {MA}$ ,所以 ${MN}//{PE}$ 
            
            又因为 ${PE}$ 在平面 ${PBC}$ 内,所以 ${MN}//$ 平面${PBC}$
            
            \item 由上一问知 ${MN}//{PE}$ ,所以 ${MN}$ 与平面${ABCD}$ 所成的角就是 ${PE}$ 与平面 ${ABCD}$ 所成的角
            
            由正棱锥的性质知${PO} = \sqrt{P{B}^{2} - O{B}^{2}} = \dfrac{{13}\sqrt{}2}{2}$ 

            ${BE} : {AD} = {BN} : {ND} = 5 : 8$ ,所以 ${BE} = \dfrac{65}{8}$

            在 $\bigtriangleup {PEB}$ 中, $\angle {PBE} = {60}^{ \circ  },{PB} = {13}$ ,
            ${BE} = \dfrac{65}{8}$ ,根据余弦定理,得 ${PE} = \dfrac{91}{8}$
            
            在 ${Rt} \bigtriangleup  {POE}$ 中, ${PO} = \dfrac{{13}\sqrt{2}}{2},{PE} = \dfrac{91}{8}$ ,所以 $\sin \angle {PEO} = \dfrac{PO}{PE} = \dfrac{4\sqrt{2}}{7}$
        \end{enumerate}
    \end{jieda}
    
    \item 如图截角四面体是一种半正八面体，可由四面体经过适当的截角，即截去四面体的四个顶点所产生的多面体.如图, 将棱长为 3 的正四面体沿棱的三等分点作平行于底面的截面得到所有棱长均为 1 的截角四面体.
    
    \begin{center}
    \includegraphics[width=0.4\textwidth]{bodyxb/src/bo_d5b4icbef24c73bio840_2_976_354_499_245_0.jpg}
    \end{center}

    \begin{enumerate}
        \item 该截角四面体的表面积
        \item 该截角四面体的体积
    \end{enumerate}

    \begin{jieda}
        \begin{enumerate}
            \item (1)依题意,该截角四面体是4个边长为1的正三角形和4个边长为1的正六边形围成,
            截角四面体中,正三角形的面积 ${S}_{1} = \dfrac{1}{2} \times  1 \times  1 \times  \dfrac{\sqrt{3}}{2} = \dfrac{\sqrt{3}}{4}$ ,
            边长为 1 的正六边形的面积 ${S}_{2} = 6 \times  \dfrac{1}{2} \times  1 \times  1 \times  \dfrac{\sqrt{3}}{2} = \dfrac{3\sqrt{3}}{2}$ ,
            所以该截角四面体的表面积为 $S = 4 \times  \dfrac{\sqrt{3}}{4} + 4 \times  \dfrac{3\sqrt{3}}{2} = 7\sqrt{3}$ .
    \item 该截角四面体是棱长为3的正四面体去掉4个角上棱长为1的正四面体而得,
    棱长为 1 的正四面体的高 $h = \sqrt{1 - {\left( \dfrac{2}{3} \times  \dfrac{\sqrt{3}}{2}\right) }^{2}} = \dfrac{\sqrt{6}}{3}$ ,棱长为 3 的正四面体的高为 ${3h} = \sqrt{6}$ ,
    则棱长为 1 的正四面体的体积 ${V}_{1} = \dfrac{1}{3} \times  \dfrac{\sqrt{3}}{4} \times  {1}^{2} \times  \dfrac{\sqrt{6}}{3} = \dfrac{\sqrt{2}}{12}$ ,
    棱长为3的正四面体的体积 ${V}_{2} = \dfrac{1}{3} \times  \dfrac{\sqrt{3}}{4} \times  {3}^{2} \times  \sqrt{6} = \dfrac{9\sqrt{2}}{4}$
    所以该截角四面体的体积为: $V = {V}_{2} - 4{V}_{1} = \dfrac{9\sqrt{2}}{4} - 4 \times  \dfrac{\sqrt{2}}{12} = \dfrac{{23}\sqrt{2}}{12}$ .
        \end{enumerate}
    \end{jieda}
    
    \item 蒙古包是蒙古族牧民居住的一种房子，建造和搬迁都很方便，适于牧业生产和游牧生活.如图 1, 一个普通的蒙古包可视为一个圆锥与一个圆柱的组合.如图 2, 已知该圆锥的高为 2 米, 圆柱的高为 3 米, 底面直径为 6 米.
    \twotu{\includegraphics[width=0.5\textwidth]{bodyxb/src/bo_d5b4icbef24c73bio840_2_707_1128_399_287_0.jpg}}{\includegraphics[width=0.5\textwidth]{bodyxb/src/bo_d5b4icbef24c73bio840_2_1124_1130_236_281_0.jpg}}

    \begin{enumerate}
        \item 求该蒙古包的侧面积
        \item 求该蒙古包的体积
    \end{enumerate}

    \begin{jieda}
        \begin{enumerate}
            \item $3\sqrt{13} \pi + 18 \pi$
            \item $33\pi$
        \end{enumerate}
    \end{jieda}
    
    \item 如图，已知正四棱锥$S-ABCD$的相邻两侧面的夹角为 ${120}^{ \circ  }$ ，它的底面边长为 $a$，求:\begin{enumerate*}
        \item 棱锥的高
        \item 斜高
        \item 侧棱长
    \end{enumerate*}
    \begin{center}
        \includegraphics[width=0.25\linewidth]{bodyxb//src/3d68d1a5748f4693e7c16193cd2d9ae2.png}
    \end{center}
    \begin{jieda}
        \begin{enumerate*}
            \item $\dfrac{a}{2}$
            \item $\dfrac{\sqrt{2}}{2}a$
            \item $\dfrac{\sqrt{3}}{2}a$
        \end{enumerate*}
    \end{jieda}

    \item 如图 (图中单位:cm) 是一种铸铁机器零件，零件下部是实心的直六棱柱 (底面是正六边形, 侧面是全等的矩形) , 上部是实心的圆柱.
        
    \begin{center}
    \includegraphics[width=0.2\textwidth]{bodyxb/src/bo_d5b4icbef24c73bio840_3_240_205_311_368_0.jpg}
    \end{center}
    \begin{enumerate}
        \item 已知铁的密度为 ${7.86}{g}/{{{cm}}}^{3}$ ，求生产一件这样的铸铁零件需要多少克铁? (结果精确到 ${0.1}{g}$ )
        \item 要给一批共 5000 个零件镀锌，若电镀这批零件每平方厘米要用锌 ${0.11}{g}$ ，求需要用锌的总量(结果精确到 ${0.01}{g}$ ).
    \end{enumerate}
    \begin{jieda}
        \begin{enumerate}
            \item 圆柱部分体积为 ${V}_{1} = \pi  \times  {3}^{2} \times  {25} = {225\pi }$ ,直六棱柱部分体积为 ${V}_{2} = 6 \times  \dfrac{\sqrt{3}}{4} \times  {12}^{2} \times  5 = {1080}\sqrt{3}$ ,则此零件的体积为 $V = {V}_{1} + {V}_{2} = {225\pi } + {1080}\sqrt{3} \approx  {2577.5}{{{cm}}}^{3}$ 
            
            又铁的密度为 ${7.86}{g}/{{{cm}}}^{3}$ ,故生产一件这样的铸铁零件需要 ${2577.5} \times  {7.86} = {20259.15} \approx  {20259.2}$ 克铁.
            
            \item 此零件的表面积为 $S = \pi  \times  {3}^{2} + {2\pi } \times  3 \times  {25} + {12} \times  5 \times  6 + 2 \times  6 \times  \dfrac{\sqrt{3}}{4} \times  {12}^{2} - \pi  \times  {3}^{2} = {150\pi } + {360} + {432}\sqrt{3} \approx  {1579.5}{{{cm}}}^{2}$ .则 5000 个零件的表面积为 ${1579.5} \times  {5000}{{{cm}}}^{2} = {7897500}{{{cm}}}^{2}$ .
            
            故需锌的质量为 ${7897500} \times  {0.11} = {868725.00}\left( {g}\right)$ .
        \end{enumerate}
    \end{jieda}
    
    \item 在四棱锥 $S - {ABCD}$ 中， ${SD}\bot$ 平面 ${ABCD}$ ，底面 ${ABCD}$ 是菱形, $E,F$ 分别为 ${SB},{AD}$ 的中点. 
            
    \begin{center}
    \includegraphics[width=0.3\textwidth]{bodyxb/src/bo_d5b4icbef24c73bio840_3_246_992_352_329_0.jpg}
    \end{center}
    \begin{enumerate}
        \item 证明: ${EF}//$ 平面 ${SCD}$
        \item 若 $\angle {BAD} = {60}^{ \circ  },{SD} = 4,{AB} = 2$ ，求三棱锥 $C - {DEF}$ 的体积.
    \end{enumerate}
    \begin{jieda}
        \begin{enumerate}
            \item 证明：如图,取 ${SC}$ 的中点 $G$ ,连接 ${DG}$ , ${EG}$ .
            $\because E$ 是 ${SB}$ 的中点, $\therefore {EG}$ 是 $\bigtriangleup {SBC}$ 的中位线,
            $\therefore {EG}//{BC},{EG} = \dfrac{1}{2}{BC}$ ,又 ${DF}//{BC},{DF} = \dfrac{1}{2}{BC}$ ,
            $\therefore {EG}//{DF},{EG} = {DF},\therefore$ 四边形 ${EGDF}$ 是平行四边形, $\therefore {EF}//{DG}$ ,
            又 ${EF} \subset$ 平面 ${SCD},{DG} \subset$ 平面 ${SCD},\therefore {EF}//$ 平面 ${SCD}$ 
            
            \item 如图,连接 ${AC}$ , ${BD}$ 交于点 $O$ ,连接 ${EO}$ ,
            $\therefore {BO} = {OD},\therefore {EO}//{SD},{EO} = \dfrac{1}{2}{SD} = 2$ ,
            又 ${SD} \bot$ 平面 ${ABCD},\therefore {EO} \bot$ 平面 ${ABCD}$ ,
            在菱形 ${ABCD}$ 中, $\angle {BAD} = {60}^{ \circ  },{AB} = 2$ ,
            $\therefore {S}_{\bigtriangleup {CDF}} = \dfrac{1}{2}{DF} \cdot  {CD}\sin {120}^{ \circ  } = \dfrac{1}{2} \times  1 \times  2 \times  \dfrac{\sqrt{3}}{2} = \dfrac{\sqrt{3}}{2}$ ,
            $\therefore {V}_{C - {DEF}} = {V}_{E - {CDF}} = \dfrac{1}{3}{EO} \cdot  {S}_{\bigtriangleup {CDF}} = \dfrac{1}{3} \times  2 \times  \dfrac{\sqrt{3}}{2} = \dfrac{\sqrt{3}}{3}$ .
        \end{enumerate}
        \begin{center}
            \includegraphics[width=0.25\linewidth]{bodyxb//src/zhuiti8.png}
        \end{center}
    \end{jieda}
\end{enumerate}

C 组：
\begin{enumerate}
    \item 已知圆锥的母线长为 4，过该圆锥顶点的平面截此圆锥所得截面面积的最大值为 8 ，则该圆锥底面半径的取值集合为\blkda{$[2\sqrt{2}, 4)$}.
    \begin{jieda}
        如图, $\bigtriangleup {ABC}$ 是圆锥的轴截面,设圆锥的底面圆半径为 $r$ .

        \begin{center}
            \includegraphics[width=0.25\linewidth]{bodyxb//src/zhuiti6.png}
        \end{center}
        
        若 $\angle {BAC} < {90}^{ \circ  }$ ,所得截面面积最大值为 ${S}_{\bigtriangleup {ABC}}$ ,则 ${S}_{\bigtriangleup {ABC}} = \dfrac{1}{2}{AB} \cdot  {AC} \cdot  \sin \angle {BAC} < \dfrac{1}{2}{AB} \cdot  {AC} = 8$ ,
        故不符合题意;
        
        若 $\angle {BAC} \geq  {90}^{ \circ  }$ ,此时所得截面面积得最大值为 ${S}_{\bigtriangleup {ABC}} = \dfrac{1}{2} \times  4 \times  4 \times  \sin {90}^{ \circ  } = 8$ ,符合题意，此时有 $\sin \angle {OAB} = \dfrac{r}{4} \geq  \sin {45}^{ \circ  }$ ,解得 $r \geq  2\sqrt{2}$ ,又 $r < 4$ ,则 $r \in  \lbrack 2\sqrt{2},4)$ .
    \end{jieda}
    \item 在三棱锥 $V - {ABC}$ 中， ${BV}\bot$ 平面 ${VAC},{VA} = 1,{AB} = {AC} = \sqrt{2},\cos \angle {VAC} = \dfrac{\sqrt{2}}{2}$ ， 点 $F$ 为棱 ${AV}$ 上一点，过点 $F$ 作三棱锥 $V - {ABC}$ 的截面，使截面平行于直线 ${VB}$ 和 ${AC}$ ，当该截面面积取得最大值时， ${CF} =$ \blkda{$\dfrac{\sqrt{5}}{2}$}.
\end{enumerate}

\section*{多面体和旋转体}
A组：
\begin{enumerate}
    \item 下列命题中正确的是\blkda{\ding{174}}.
    \begin{dingautolist}{172}
        \item 圆柱是将矩形旋转一周所得的几何体
        \item 平行于圆锥某一母线的截面是等腰三角形
        \item 过圆锥的轴的截面是等腰三角形
        \item 过直角三角形的一边为旋转轴, 旋转所得的几何体是圆锥
    \end{dingautolist}
    
    \item 三棱锥的三条侧棱两两垂直，且这三条侧棱长分别是 1、2、3，那么该棱锥的体积为\blkda{1}.
    \begin{jieda}
        $V = \dfrac{1}{3} \times \dfrac{1}{2} \times 1 \times 2 \times 3 = 1$
    \end{jieda}
    
    \item 已知正三棱锥底面边长是 4，侧棱长为 3，则其体积为\blkda{$\dfrac{4\sqrt{11}}{3}$}.
    \begin{jieda}
        底面顶点到中心距离为$\dfrac{4\sqrt{3}}{3}$，三棱锥的高为$\dfrac{\sqrt{33}}{3}$，体积$V = \dfrac{1}{3} \times \dfrac{\sqrt{3}}{4} \times 16 \times \dfrac{\sqrt{33}}{3} = \dfrac{4\sqrt{11}}{3}$
    \end{jieda}
    
    \item 在三棱台 ${ABC} - {A}_{1}{B}_{1}{C}_{1}$ 中，三棱锥 $B - {A}_{1}{B}_{1}{C}_{1}$ 体积为 1，三棱锥 ${C}_{1} - {ABC}$ 体积为 4，则该三棱台体积为\blkda{7}.
    \begin{jieda}
        $V = \dfrac{1}{3} \times S \times h = 1$，$V' = \dfrac{1}{3} \times S' \times h = 4$， $V_{\text{台}} = \dfrac{1}{3}(S+S'+\sqrt{SS'}) \times h = 1+4+\sqrt{1\times4} = 7$
    \end{jieda}
    
    \item 若用长为 ${10}{{cm}}$ ，宽为 $6{{cm}}$ 的长方形卷成一个圆柱，则此圆柱的体积为\blkda[3]{$\dfrac{150}{\pi} cm^3$或$\dfrac{90}{\pi}cm^3$}.
    
    \item 若把半径为 $3{{cm}}$ 、圆心角为 $\dfrac{2}{3}\pi$ 的扇形卷成一个圆锥，则此圆锥的体积为\blkda{$\dfrac{2\sqrt{2}\pi}{3}$}.
    
    \item 如果圆锥的底面半径为 $\sqrt{2}$ ，高为 2 ，那么此圆锥的侧面积是\blkda{$2\sqrt{3}\pi$}.
    \begin{jieda}
        $l = \sqrt{6}$，$S = \pi rl = 2\sqrt{3}\pi$
    \end{jieda}

\end{enumerate}
B 组：
\begin{enumerate}
    \item 设一个圆柱的底面半径为 $5{{cm}}$ ，高 $8{{cm}}$ ，平行于轴的平面截圆柱所得的截面面积 ${48c}{m}^{2}$ ，则此截面与轴的距离为\blkda{$4cm$}.

    \item 经过高为 ${10cm}$ 的圆锥的顶点,且与底面成 ${45}^{ \circ  }$ 二面角的平面把圆锥底面圆的周长截取 $\dfrac{1}{4}$ ，则此截面面积为\blkda{$100\sqrt{2}$}平方厘米.
    \begin{jieda}
        设二面角棱为$AB$，底面截取的扇形$OAB$圆心角为$\dfrac{\pi}{2}$，取$AB$中点$M$，则$OM = 10cm$，故底面半径为$10\sqrt{2} cm$
        \begin{enumerate}
            \item 方法1

            因此母线长为$10\sqrt{3} cm$，因此截面为$10\sqrt{3}cm, 10\sqrt{3}cm, 20cm$的等腰三角形，根据余弦定理可知其顶角余弦值为$\dfrac{1}{3}$，所以其正弦值为$\dfrac{2\sqrt{2}}{3}$，所以三角形面积为$100\sqrt{2} cm^2$
            
            \item 方法2

            截面在底面投影面积$S' = 100cm^2$
            $S = \dfrac{S'}{cos\ang{45}} = 100\sqrt{2} cm^2$
        \end{enumerate}
    \end{jieda}

    \item 设圆锥的母线长等于高的 2 倍，则过它顶点且有最大截面积的截面与底面所成的角是\blkda{$\dfrac{\pi}{4}$}.
    \begin{jieda}
        $r:h:l = \sqrt{3}:1:2$，最大截面面积$S = \dfrac{l^2}{2} = 2h^2$，截面和底面交线长度为$2\sqrt{2}h$，故截面在底面投影面积$S' = \dfrac{1}{2} \times 2\sqrt{2}h \times h = \sqrt{2}h^2$，故$cos\theta = \dfrac{\sqrt{2}}{2}$，$\theta = \dfrac{\pi}{4}$
    \end{jieda}

    \item 在侧棱长为 2 的正三棱锥中, 若其底面周长为 9 , 则该正三棱锥的体积是\blkda{$\dfrac{3\sqrt{3}}{4}$}.
    \begin{jieda}
        高$h = 1$，底面面积$S = \dfrac{\sqrt{3}}{4} \times 3^2 = \dfrac{9\sqrt{3}}{4}$
        
        $V = \dfrac{1}{3} \times \dfrac{9\sqrt{3}}{4} \times 1 = \dfrac{3\sqrt{3}}{4}$
    \end{jieda}

    \item 若正方形 ABCD 边长为 1, E、F 分别为 BC、CD 中点, 沿 AE、EF、AF 折成一个三棱锥，则这个三棱锥的体积是\blkda{$\dfrac{1}{24}$}.
    \begin{jieda}
        设$B,C,D$折叠后重合在$P$，则$P$为墙角$V = \dfrac{1}{3} \times \left(\dfrac{1}{2} \times \dfrac{1}{2} \times \dfrac{1}{2}\right) \times 1 = \dfrac{1}{24}$
    \end{jieda}

    \item 设正四棱柱的一条对角线长为 $3{{cm}}$ ，它的全面积是 ${16}{{{cm}}}^{2}$ ，则它的体积是\blkda{$4$或$\dfrac{112}{27}$}.
    \begin{jieda}
        设底面边长为$a$厘米，高为$h$厘米，则$\left\{\begin{array}{l}
             2a^2+4ah = 16 \\
             a^2+a^2+h^2 = 9
        \end{array}\right.$，解得$a = 2, h = 1$或$a = \dfrac{4}{3}, h = \dfrac{7}{3}$，因此$V = a^2h = 4$或$\dfrac{112}{27}$
    \end{jieda}

    \item 若底面为平行四边形的直棱柱的侧棱长为 $5{{cm}}$ ，底面边长分别为 $6{{cm}}$ 和 $8{{cm}}$ ， 且两边长的夹角为 ${30}^{ \circ  }$ ，则此直棱柱的全面积为\blkda{$188cm^2$}
    \begin{jieda}
        $S_{\text{底}} = 24cm^2$，$S_{\text{侧}} = 140cm^2$，$S_{\text{表}} = 188cm^2$
    \end{jieda}

    \item 若正三棱锥的一个侧面的面积与底面的面积之比等于 2:3 , 则这个三棱锥的侧面和底面所成的二面角是\blkda{$\dfrac{\pi}{3}$}.
    \begin{jieda}
        $cos\theta = \dfrac{S'}{S} = \dfrac{1}{2}$，$\theta = \dfrac{\pi}{3}$
    \end{jieda}

    \item 如图，有一堆规格相同的铁制(铁的密度是 ${7.8}{{g/c{m}^{3}}}$ )六角螺帽共重 5.8kg，已知底面是正六边形，边长为 12mm，内孔直径为 10mm，高为 10mm，问这堆螺帽大约有\blkda{252}个
    
    \begin{center}
    \includegraphics[width=0.2\textwidth]{bodyxb/src/bo_d5fmdlf7aajc7382at8g_1_1193_544_236_174_0.jpg}
    \end{center}


    \item 圆柱 $O{O}_{1}$ 的底面半径为 $R$ ，上底半径 ${OA}$ 与下底半径 ${O}_{1}{B}_{1}$ 的夹角为 $\theta$ ，则 $O{O}_{1}$ 与 $A{B}_{1}$ 的距离为\blkda[3]{$R\cos \dfrac{\theta }{2}$或$R\sin \dfrac{\theta }{2}$}
    

    \item 若一个正四棱锥的中截面（过各侧棱中点的截面）面积是 $Q$ ,则它的底面边长是\dotfill{(\kh{D})}
    \xx{$\dfrac{Q}{4}$}{$\dfrac{\sqrt{Q}}{2}$}{$\sqrt{Q}$}{$2\sqrt{Q}$}
    \begin{jieda}
        $S_{\text{底}} = 4Q$，故底面边长为$2\sqrt{Q}$
    \end{jieda}
    
    \item 若三棱锥 $A-BCD$ 的侧面 $ABC$ 内一点 $P$ 到底面 $BCD$ 的距离与到棱 $AB$ 的距离相等，则动点 ${P}$ 的轨迹是\dotfill{(\kh{A})}
    \xx{线段}{抛物线的一部分}{双曲线的一部分}{一段圆弧}
    \begin{jieda}
        $P$ 到底面 $BCD$ 的距离可转化为$k$倍的$P$到$BC$的距离，故$P$到$AB$和$BC$的距离之比为定值，因此其轨迹为线段.
    \end{jieda}
    
    
    \item 两相同的正四棱锥组成如图所示的几何体, 可放棱长为 1 的正方体内，使正四棱锥的底面 ABCD 与正方体的某一个平面平行，且各顶点均在正方体的面上，则这样的几何体体积的可能值有\dotfill{(\kh{D})}
    \xx{1 个}{2 个}{3 个}{无穷多个}

    \twotu{\includegraphics[width=0.4\textwidth]{bodyxb/src/bo_d5fmdlf7aajc7382at8g_1_1152_1477_180_202_0.jpg}}{\includegraphics[width=0.4\textwidth]{bodyxb/src/bo_d5fmdlf7aajc7382at8g_1_1356_1450_242_236_0.jpg}}
    \begin{jieda}
        问题转化为为正方形有多少个内接正方形
    \end{jieda}
    
    
    \item 已知在平行六面体 ${ABCD} - {A}_{1}{B}_{1}{C}_{1}{D}_{1}$ 中，底面 ${ABCD}$ 为矩形， ${AB} = 8$ ， ${AD} = {10}$ ， ${{{AA}}}_{1} = 6$ ，且 $\angle {{{AA}}}_{1}{{B}}_{1} = \angle {{{AA}}}_{1}{{D}}_{1} = \dfrac{\pi }{3}$ ，求此平行六面体的体积
    \begin{jieda}
        侧棱和底面的线面角为$\dfrac{\pi}{4}$，所以平行六面体的高为$3\sqrt{2}$，因此体积$V = 8 \times 10 \times 3\sqrt{2} = 240\sqrt{2}$
    \end{jieda}
    
    \item 已知有一个高是 $8{{cm}}$ 的圆柱,用一个平行于轴且和轴相距 $2{{cm}}$ 的平面去截它, 此截面在圆柱底面上截得的弧所对的圆心角是 ${60}^{ \circ  }$ ,求此截面的面积
    \begin{jieda}
        截面和底面的交线长为$\dfrac{4\sqrt{3}}{3} cm$，截面面积为$\dfrac{32\sqrt{3}}{3} cm^2$.
    \end{jieda}
    
    \item 已知圆柱的轴为 $O{O}_{1}$ ， $A$ ， $B$ 分别为上、下底面圆周上的点，若 $O{O}_{1} = 4$ ， ${OA} = 3$ ， 且 ${OA}$ 与 ${O}_{1}B$ 成 ${60}^{ \circ  }$ 角
    \begin{enumerate*}
        \item 求 $O{O}_{1}$ 与 ${AB}$ 所成角的正切值
        \item 求过 ${AB}$ 且与 $O{O}_{1}$ 平行的平面和 $O{O}_{1}$ 的距离
    \end{enumerate*}
    \begin{jieda}
        \begin{enumerate}
            \item 如图,过点 $A$ 作圆柱的母线 $A{A}_{1}$ ,交上底面圆周于点 ${A}_{1}$ ,连接 ${O}_{1}{A}_{1}$ ,连接 ${A}_{1}B$ ,则 ${O}_{1}{A}_{1}//{OA}$ ,且 $A{A}_{1}//O{O}_{1}$ , $A{A}_{1} = O{O}_{1} = 4$
            \begin{center}
                \includegraphics[width=0.25\linewidth]{bodyxb//src/yuanzhu1.png}
            \end{center}
            
            $\angle {A}_{1}{AB}$ (或其补角)就是 $O{O}_{1}$ 与 ${AB}$ 所成的角
            
            $\angle {A}_{1}{O}_{1}B$ 就是 ${OA}$ 与 ${O}_{1}B$ 所成的角,即 $\angle {A}_{1}{O}_{1}B = {60}^{ \circ  }$或$\ang{120}$ 

            故${A}_{1}B = 3$或$3\sqrt{3}$
            
            $\because A{A}_{1} \bot$ 平面 ${O}_{1}{A}_{1}B,{A}_{1}B \subset$ 平面 ${O}_{1}{A}_{1}B$ ,
            $\therefore A{A}_{1}\bot {A}_{1}B$ 
            
            在 ${Rt} \bigtriangleup  A{A}_{1}B$ 中, $\tan \angle {A}_{1}{AB} = \dfrac{{A}_{1}B}{A{A}_{1}} = \dfrac{3}{4}$或$\dfrac{3\sqrt{3}}{4}$
            
            $\therefore O{O}_{1}$ 与 ${AB}$ 所成角的正切值为 $\dfrac{3}{4}$或$\dfrac{3\sqrt{3}}{4}$
            \item 由 ${O}_{1}O//A{A}_{1}$ ,得 $O{O}_{1}//$ 平面 $A{A}_{1}B$ . 取 ${A}_{1}B$ 的中点 $C$ ,连接 ${O}_{1}C$ ,则 ${O}_{1}C \bot  {A}_{1}B$ .
            
            又 $\because {O}_{1}C \bot  A{A}_{1},\therefore {O}_{1}C \bot$ 平面 $A{A}_{1}B$ ,故 ${O}_{1}C$ 就是 $O{O}_{1}$ 与平面 $A{A}_{1}B$ 间的距离.
            
            容易算得 ${O}_{1}C = \dfrac{3\sqrt{3}}{2}$ 或 $\dfrac{3}{2}$ .
        \end{enumerate}
    \end{jieda}
    
    \item 如图，梯形 ${ABCD}$ 满足 ${AB}//{CD},\angle {ABC} = {90}^{ \circ  }$ ，且 ${AB} = {2\sqrt{3}},{BC} = 1$ ， $\angle {BAD} = {30}^{ \circ  }$ ,现将梯形 ${ABCD}$ 绕 ${AB}$ 所在的直线旋转一周,所得几何体记作 $\Omega$ .
    
    \begin{center}
    \includegraphics[width=0.1\textwidth]{bodyxb/src/bo_d5fmdlf7aajc7382at8g_2_1288_1326_169_478_0.jpg}
    \end{center}
    
    \begin{enumerate*}
        \item 求 $\Omega$ 的体积 $V$
        \item 求 $\Omega$ 的表面积 $S$
    \end{enumerate*}
    \begin{jieda}
        \begin{enumerate}
            \item $\dfrac{4\sqrt{3}}{3}\pi$
            \item $(3+2\sqrt{3})\pi$
        \end{enumerate}
    \end{jieda}
\end{enumerate}
\section*{球（1）}
A组：
\begin{enumerate}
    \item 过球面上任意两点作大圆，可作大圆\dotfill{(\kh{C})}
    \xx{1 个}{无穷多个}{1 或无穷多个}{0 个}
    
    \item 下列说法中不正确的是\dotfill{(\kh{C})}
    \xx{经过圆锥的轴的截面是一个等腰三角形}{矩形绕一边旋转一周得到圆柱}{圆柱的截面一定是矩形面}{球的截面一定是圆面}
    
    \item 已知球的大圆的周长是 $C$ ，则此球的表面积是\blkda{$\dfrac{C^2}{\pi}$}.
    \begin{jieda}
        $C = 2\pi r$，$S = 4\pi r^2 = \dfrac{C^2}{\pi}$
    \end{jieda}
    
    \item 两球面积之差为 ${48\pi }$ ，大圆周长和为 ${12\pi }$ ，则两球的半径为\dotfill{(\kh{A})}
    \xx{$2,4$}{${2\pi },{4\pi }$}{$6,4$}{${6\pi },{4\pi }$}
    \begin{jieda}
        $\left\{\begin{array}{l}
             4\pi(R^2 - r^2) = 48\pi  \\
             2\pi (R+r) = 12\pi
        \end{array}\right.$
    \end{jieda}
\end{enumerate}
B组：
\begin{enumerate}
    \item 在半径为$10cm$ 的球内，有一个截面距球心$6cm$，则该截面面积等于\blkda{$64\pi  cm^2$}.
    \begin{jieda}
        $r = \sqrt{R^2-d^2} = 8cm$，$S = \pi r^2 = 64\pi cm^2$
    \end{jieda}

    \item 已知半径为 $13$ 的球的两个平行截面的周长为 $10\pi$ 和 $24\pi$，则这两个截面间的距离为\blkda{$7$或$17$}.
    \begin{jieda}
        $r_1 = 5, r_2 = 12$，$d_1 = 12, d_2 = 5$
    \end{jieda}

    \item 已知过球的某条半径的中点 ${O}_{1}$ 作一个垂直于这条半径的平面，求所得截面面积和球大圆面积之比为\blkda{$\dfrac{3}{4}$}.
    \begin{jieda}
        $\dfrac{r}{R} = \dfrac{\sqrt{3}}{2}$，$\dfrac{S}{S'} = \dfrac{3}{4}$
    \end{jieda}

    \item 若球的表面积扩大为原来的 2 倍，则体积是原来的\blkda{$(\sqrt{2})^3$}倍.
    \begin{jieda}
        $\dfrac{r_1}{r_2} = \sqrt{\dfrac{S_1}{S_2}} = \sqrt{2}$，$\dfrac{V_1}{V_2} = \left(\dfrac{r_1}{r_2}\right)^3 = (\sqrt{2})^3$
    \end{jieda}

    \item 一球的体积和表面积在数值上相等，则该球的半径数值为\blkda{3}.
    \begin{jieda}
        $\dfrac{4}{3}\pi R^3 = 4\pi R^2$，$R = 3$
    \end{jieda}

    \item 球面上三点 ${A}\text{ 、 }{B}\text{ 、 }{C},{{AB}} = {10},\angle {C} = {30}^{ \circ  }$ ,球的半径为 20，则球心到平面的距离为\blkda{$10\sqrt{3}$}.
    \begin{jieda}
        $2r = \dfrac{10}{sin\ang{30}} = 20$，故$r = 10$，因此$d = \sqrt{R^2-r^2} = 10\sqrt{3}$
    \end{jieda}

    \item 体积相等的球、正四面体和正方体，则它们的表面积间的大小关系为\dotfill{(\kh{B})}
    \xx{ ${S}_{\text{ 球 }} < {S}_{\text{ 正四面体 }} < {S}_{\text{ 正方体 }}$}{${S}_{\text{ 球 }} < {S}_{\text{ 正方体 }} < {S}_{\text{ 正四面体 }}$}{${S}_{\text{ 正四面体 }} < {S}_{\text{ 球 }} < {S}_{\text{ 正方体 }}$}{${S}_{\text{ 正方体 }} < {S}_{\text{ 球 }} < {S}_{\text{ 正四面体 }}$}
    
    \item 若一平面截球得到直径是$6cm$ 的圆面，球心到这个平面的距离是 $4cm$，则该球的体积是\blkda{$\dfrac{500\pi}{3}$}.
    \begin{jieda}
        $r = 3, R = 5, V = \dfrac{500\pi}{3}$
    \end{jieda}

    \item 圆柱的底面直径与高都等于球的直径，则球的表面积与圆柱的侧面积的比值为\dotfill{(\kh{A})}
    \xx{$1 : 1$}{$1 : 2$}{$2 : 1$}{$2 : 3$}
    \begin{jieda}
        $S_{\text{侧}} = 2\pi r \times (2r) = 4\pi r^2$， $S_{\text{表}} = 4\pi r^2$
    \end{jieda}

    \item 连接球面上两点的线段称为球的弦.半径为 4 的球的两条弦 ${AB},{CD}$ 的长度分别等于 $2\sqrt{7},4\sqrt{3},M,N$ 分别为 ${AB},{CD}$ 的中点,每条弦的两端都在球面上运动,有下列四个命题
    \begin{dingautolist}{172}
        \item 弦 ${AB},{CD}$ 可能相交于点 $M$ 
        \item 弦 ${AB},{CD}$ 可能相交于点 $N$
        \item ${MN}$ 的最大值为5
        \item ${MN}$ 的最小值为1.
    \end{dingautolist}
    其中真命题的个数是\blkda{3}个.
    \begin{jieda}
        $r_1 = \sqrt{7}, d_1 = 3$，$r_2 = 2\sqrt{3}, d_2 = 2$

        若两弦交于 $N$ ,则 ${OM}\bot {MN}$ , ${{Rt}}\triangle {OMN}$ 中,有 ${OM} < {ON}$ ,矛盾，故\ding{173}错.
    \end{jieda}
    
    \item 已知 $P$ 是半径为 $R$ 的球面上任意一点， ${PA},{PB},{PC}$ 是两两互相垂直的三条弦. 求证: $P{A}^{2} + P{B}^{2} + P{C}^{2}$ 是一个定值.
    \begin{jieda}
    设球心为 $O$ ,过 ${PA}\text{、}{PB}$ 作平面,截球得截面圆,设截面圆的圆心为 ${O}_{1}$ ,连接 $P{O}_{1}$并延长交圆${O}_{1}$ 于点 $D$ ,则 ${APBD}$ 为矩形$C$ 、 $O$ 、 $D$ 三点共线
    
    $P{A}^{2} + P{B}^{2} + P{C}^{2} = A{B}^{2} + P{C}^{2} = P{D}^{2} + P{C}^{2} = C{D}^{2}$ 为定值
    
    \ps 题干条件即三条弦${PA},{PB},{PC}$为邻边的长方体的外接球半径为$R$，故$P{A}^{2} + P{B}^{2} + P{C}^{2} = 4R^2$
    \end{jieda}

    \item 已知 $A,B,C$ 为半径为 $2\sqrt{3}$ 的球面上三点， ${AB} = 2,{BC} = 4$ ， $\angle {ABC} = {60}^{ \circ  }$ ， $O$ 为球心,求直线 ${OA}$ 与截面 ${ABC}$ 所成角.
    \begin{jieda}
        易知$\triangle ABC$为直角三角形，取斜边$BC$中点$M$则$OM = \sqrt{12-4} = 2\sqrt{3}$

        设直线 ${OA}$ 与截面 ${ABC}$ 所成角为$\theta$，则$sin\theta = \dfrac{OM}{OA} = \dfrac{2\sqrt{2}}{2\sqrt{3}} = \dfrac{\sqrt{6}}{3}$，故直线 ${OA}$ 与截面 ${ABC}$ 所成角为$arcsin\left(\dfrac{\sqrt{6}}{3}\right)$
    \end{jieda}
    
    \item 已知球面上过 $A,B,C$ 三点的截面和球心的距离等于球半径的一半,且 ${AB} = {BC} = {CA} = 2$ ，求此球面的面积.
    \begin{jieda}
        设正三角形$ABC$的中心为$O_1$，球心为$O$，则$OO_1 = \dfrac{OA}{2}$，$O_1A = \sqrt{OO_1} = \dfrac{2}{\sqrt{3}}$，故$R = OA = \dfrac{4}{3}$，$S = 4\pi R^2 = \dfrac{64\pi}{9}$
    \end{jieda}
\end{enumerate}

补充练习：
\begin{enumerate}
    \item 若圆柱的底面半径为 13，平行于圆柱的轴的截面面积是 48，与轴的距离是 5， 则母线的长为\blkda{2}.

    \item 过高为 ${10}{{cm}}$ 的圆锥顶点作与底面成 ${60}^{ \circ  }$ 角的截面，此截面将圆周底面周长截成1:2的部分，则此截面的面积等于\blkda{$\dfrac{200\sqrt{3}}{3}cm^2$}.

    \item 若一个等边圆柱(轴截面是正方形的圆柱)的轴截面的面积是 $S$ ，则它的一个底面的面积是\blkda{$\dfrac{\pi S}{4}$}.

    \item 若过圆锥顶点的截面的三角形面积的取值范围为 $(0,4\sqrt{3}\rbrack$ ，该圆锥的母线长为 4 ，则该圆锥顶角的最大值为\blkda{$\dfrac{\pi}{3}$}.

    \item 如图(1)所示，一个正三棱柱形容器，高为 ${2a}$ ，内装水若干，将容器放倒，把一个侧面作为底面，如图(2)所示， 这时水面恰为中截面，则图(1)中水面的高度是\blkda{$\dfrac{3a}{2}$}.
    \twotu{\includegraphics[width=0.3\textwidth]{bodyxb/src/bo_d5gu7tbef24c73bktr7g_2_219_822_181_235_0.jpg}}{\includegraphics[width=0.5\textwidth]{bodyxb/src/bo_d5gu7tbef24c73bktr7g_2_416_886_228_147_0.jpg}}

    \item 设正多面体的每个面都是正 $n$ 边形，以每个顶点为端点的棱有 $m$ 条，棱数是 $E$，面数是 ${F}$ ，则它们之间的关系不正确的是\dotfill{(\kh{D})}
    \xx{${{nF}} = 2{E}$}{${{mV}} = 2{E}$}{$V+F=E+2$}{${{mF}} = 2{E}$}
    
    \item 如图所示的几何体是一棱长为 $4{{cm}}$ 的正方体，若在它的各个面的中心位置上， 各打一个直径为 $2{{cm}}$ 、深为 $1{{cm}}$ 的圆柱形的孔，求打孔后几何体的表面积是多少? ($\pi$取 3.14)
    
    \begin{center}
    \includegraphics[width=0.2\textwidth]{bodyxb/src/bo_d5gu7tbef24c73bktr7g_2_1127_1605_277_264_0.jpg}
    \end{center}
    \begin{jieda}
        ${96} + {12\pi } \approx  {133.68}$
    \end{jieda}


    \item 如图，在三棱锥$P-ABC$ 中，$PA\bot BC$，$PA=BC=1$, $PA$、$BC$的公垂线段$EF=h$， 求三棱锥 $P-ABC$ 的体积.
    
    \begin{center}
    \includegraphics[width=0.2\textwidth]{bodyxb/src/bo_d5gu7tbef24c73bktr7g_3_1052_303_325_270_0.jpg}
    \end{center}
    \begin{jieda}
         $\dfrac{h}{6}$
    \end{jieda}

    \item 设 $a > 2,m \in  {R},k,t$ 是不全为零的实数，椭圆 $\Gamma  : \dfrac{{x}^{2}}{{a}^{2}} + \dfrac{{y}^{2}}{4} = 1,{F}_{1}\text{ 、 }{F}_{2}$ 分别为 $\Gamma$ 的左、右两个焦点,直线 $l : {kx} + {ty} + m = 0$ 与 $\Gamma$ 交于 $A,B$ 两点， $O$ 为坐标原点.
    \begin{enumerate}
        \item 若 $k = 0,t + m = 0$ ，且四边形 ${AB}{F}_{1}{F}_{2}$ 为矩形，求 $\Gamma$ 的离心率
        \item 若 $t = 0$ ，且 $\bigtriangleup  {F}_{1}{AB}$ 的周长的最大值为 12，求 $\Gamma$ 的方程
        \item 若 $a = 4,t =  - 1,\lambda  > 0$ ，直线 ${OA}$ 与 ${OB}$ 的斜率之积为 $- \dfrac{1}{4}$ ， $C$ 为 ${OA}$ 上的一点， 且 $\vv{OA} = \vv{AC}$ ,直线 ${BC}$ 与椭圆 $\Gamma$ 交于点 $D$ ,且 $\vv{BD} = \lambda \vv{DC}$ ,求 $\lambda$ 的值以及 $\bigtriangleup {ABD}$ 的面积 ${S}$ 的值
    \end{enumerate}
    \begin{jieda}
    \begin{enumerate}
        \item 由 $k = 0,t + m = 0$ 得, $l : y = 1$，将 $y = 1$ 代入椭圆方程得 $\dfrac{{x}^{2}}{{a}^{2}} + \dfrac{1}{4} = 1$ ,解得 $x =  \pm  \dfrac{\sqrt{3}}{2}a$ ,则 $\left| {AB}\right|  = \sqrt{3}a$
        
        又四边形 ${AB}{F}_{1}{F}_{2}$ 是矩形,则 $\sqrt{3}a = {2c}$ ,即离心率 $e = \dfrac{c}{a} = \dfrac{\sqrt{3}}{2}$ .
        \item 由 $t = 0$ 得， $l : x =  - \dfrac{m}{k}$ ，即 ${AB}\bot x$ 轴，则 $\left| {A{F}_{1}}\right|  + \left| {B{F}_{1}}\right|  + \left| {AB}\right|  \leq  \left| {A{F}_{1}}\right|  + \left| {B{F}_{1}}\right|  + \left| {A{F}_{2}}\right|  + \left| {B{F}_{2}}\right|  = {4a}$
        
        当且仅当 ${AB}$ 过右焦点 ${F}_{2}$ 时等号成立,即 $\bigtriangleup  {F}_{1}{AB}$ 的周长的最大值为 ${4a}$，即 ${4a} = {12}$ ,即 $a = 3$ , 则 $\Gamma$ 方程为 $\dfrac{{x}^{2}}{9} + \dfrac{{y}^{2}}{4} = 1$ .
        
        \begin{center}
        \includegraphics[width=0.3\textwidth]{bodyxb/src/bo_d5gv9kref24c73bkubog_1_243_609_378_305_0.jpg}
        \end{center}

        \item 设$A(4cos\alpha, 2sin\alpha)$，$B(4cos\beta, 2sin\beta)$，
        
        则有$C(8cos\alpha, 4sin\alpha)$，$D\left(\dfrac{\lambda(8cos\alpha)+(4cos\beta)}{1+\lambda}, \dfrac{\lambda(4sin\alpha)+(2sin\beta)}{1+\lambda}\right)$

        $k_{OA} = \dfrac{1}{2}tan\alpha$，$k_{OB} = \dfrac{1}{2}tan\beta$，故$tan\alpha \cdot tan\beta = -1$

        将$D$的坐标代入椭圆方程$\left[\dfrac{\lambda(8cos\alpha)+(4cos\beta)}{1+\lambda}\right]^2+4\left[\dfrac{\lambda(4sin\alpha)+(2sin\beta)}{1+\lambda}\right]^2 = 16$

        即$[\lambda(2cos\alpha)+cos\beta]^2+[\lambda(2sin\alpha)+sin\beta]^2 = (1+\lambda)^2$

        即$4\lambda^2 + 4cos(\alpha - \beta)\lambda = \lambda^2+2\lambda$，代入$tan\alpha \cdot tan\beta = -1$即$3\lambda^2 = 2\lambda$

        故$\lambda = \dfrac{2}{3}$或$0$（舍）

        $S_{\triangle ABD} = \dfrac{2}{5}S_{\triangle ABC} = \dfrac{2}{5}S_{\triangle ABO}$

        $S_{\triangle ABO} = 4\left|sin(\alpha - \beta)\right| = 4$，因此$S_{\triangle ABD} = \dfrac{8}{5}$

        
        \begin{center}
        \includegraphics[width=0.3\textwidth]{bodyxb/src/bo_d5gv9kref24c73bkubog_2_241_1243_424_254_0.jpg}
        \end{center}
    \end{enumerate}
    \end{jieda}
\end{enumerate}

\section*{球（2）}
\begin{enumerate}
    \item 将一个半径为 $R$ 的木球削成一个尽可能大的正方体，则此正方体的体积是\blkda{$\dfrac{8\sqrt{3}}{9}R^3$}.
    \begin{jieda}
        正方体的外接球半径为$R$，设正方体的棱长为$a$，则$2R = \sqrt{3}a$，$V = a^3 = \dfrac{8\sqrt{3}}{9}R^3$
    \end{jieda}

    \item 把三个半径为 $R$ 的铁球熔铸成一个底面半径为 $R$ 的圆柱，不计损耗，则圆柱的高为\blkda{$4R$}.

    \item 正方体的全面积为 ${a}^{2}$ ，它的顶点都在球面，这个球的表面积是\blkda{$\dfrac{\pi a^2}{2}$}.
    \begin{jieda}
        正方体棱长$\dfrac{a}{\sqrt{6}}$，求的直径$\dfrac{a}{\sqrt{2}}$，求的表面积$4\pi \times \left(\dfrac{\sqrt{2}a}{4}\right)^2 = \dfrac{\pi a^2}{2}$
    \end{jieda}

    \item 已知半径为 $R$ 的半球内有一内接圆柱，此圆柱的全面积的最大值是\blkda[3]{$\left( {1 + \sqrt{2}}\right) \pi {R}^{2}$}.
    \begin{jieda}
        如图设圆柱的底面半径为 $r$ ,高为 $h$ ,则 $h = R\sin \theta ,r = R\cos \theta$ 

        \begin{center}
            \includegraphics[width=0.25\linewidth]{bodyxb//src/qiu1.png}
        \end{center}
        
        设圆柱的表面积为 $S$ ,则
        $S = {2\pi }{r}^{2} + {2\pi rh} = {2\pi }{R}^{2}{\cos }^{2}\theta  + {2\pi }{R}^{2}\sin \theta \cos \theta$
        $= \pi {R}^{2}\left( {\sin {2\theta } + \cos {2\theta } + 1}\right)$
        $= \pi {R}^{2}\left\lbrack  {\sqrt{2}\sin \left( {{2\theta } + \dfrac{\pi }{4}}\right)  + 1}\right\rbrack$ ,
        $\therefore$ 当 ${2\theta } + \dfrac{\pi }{4} = \dfrac{\pi }{2}$ ,即 $\theta  = \dfrac{\pi }{8}$ 时, $S$ 有最大值 $\left( {1 + \sqrt{2}}\right) \pi {R}^{2}$ .
    \end{jieda}

    \item 已知 $P$、$A$、$B$、$C$ 是球 $O$ 的面上的四个点，$PA$、$PB$、$PC$ 两两垂直，且 $PA=PB=PC=1$， 则球 $O$ 的体积\blkda{$\dfrac{\sqrt{3}}{2}\pi$}，球 $O$ 的表面积是\blkda{$3\pi$}.
    \begin{jieda}
        $R = \dfrac{\sqrt{3}}{2}$，$V = \dfrac{4\pi}{3}\left(\dfrac{\sqrt{3}}{2}\right)^3 = \dfrac{\sqrt{3}}{2}\pi$，$S = 3\pi$
    \end{jieda}

    \item 已知有四个半径为 $R$ 的球，每个球都和其它三球相切，和这四个球都相切的球的半径是\blkda[5]{$\left(1+\dfrac{\sqrt{6}}{2}\right)R$或$\left(\dfrac{\sqrt{6}}{2}-1\right)R$}.
    \begin{jieda}
        四个球心组成棱长为$2R$的正四面体，则正四面体中心到顶点的距离为$\dfrac{\sqrt{6}}{2}R$，故所求的半径为$\left(1+\dfrac{\sqrt{6}}{2}\right)R$或$\left(\dfrac{\sqrt{6}}{2}-1\right)R$
    \end{jieda}

    \item 半径为 $R$ 的 5 个球中的四个放在桌面上，使得任意一个球仅和其中的两个球相切，第五个球放在这四个球上，使它与这四个球都相切，则第五个球球心到桌面的距离为\dotfill{(\kh{A})}
    \xx{$\left( {1 + \sqrt{2}}\right) R$}{$\sqrt{2}R$}{$2\sqrt{2}R$}{$\left( {1 + 2\sqrt{2}}\right) R$}
    
    \item 取一个底面半径和高都为 $R$ 的圆柱,从圆柱中挖去一个以圆柱的上底面为底面,下底面圆心为顶点的圆锥,把所得的几何体与一个半径为 $R$ 的半球放在同一水平面 $a$ 上. 用一平行于平面 $a$ 的平面去截这两个几何体，截面分别为圆面和圆环面. 设截面面积分别为 ${S}_\text{圆}$ 和 ${{S}}_{\text{ 圆环 }}$ ,那么 \dotfill{(\kh{C})}
    \xx{${S}_{\text{ 圆 }} > {S}_{\text{ 圆环 }}$}{${S}_{\text{ 圆 }} < {S}_{\text{ 圆环 }}$}{${S}_{\text{ 圆 }} = {S}_{\text{ 圆环 }}$}{不确定}
    
    \item 如图，有一个水平放置的透明无盖的正方体容器，容器高8cm，将一个球放在容器口，再向容器内注水，当球面恰好接触水面时测得水深为 6cm，如果不计容器的厚度，则球的体积为\dotfill{(\kh{A})}
    
    \begin{center}
    \includegraphics[width=0.2\textwidth]{bodyxb/src/bo_d5gv3rn7aajc7382nmmg_0_227_1612_235_309_0.jpg}
    \end{center}

    \xx{$\dfrac{500\pi }{3}c{m}^{3}$}{$\dfrac{866\pi }{3}c{m}^{3}$}{$\dfrac{1372\pi }{3}c{m}^{3}$}{$\dfrac{2048\pi }{3}c{m}^{3}$}
    
    \item 一个倒圆锥形容器，它的轴截面是正三角形，在容器内注入水，并放入一个半径为 $r$ 的铁球, 这时水面恰好和球面相切. 问将球从圆锥内取出后, 圆锥内水平面的高是多少?
    \begin{jieda}
        如图作轴截面
        
        \begin{center}
            \includegraphics[width=0.25\linewidth]{bodyxb//src/qiu2.png}
        \end{center}
        设球未取出时水面高 ${PC} = h$ ,球取出后,水面高 ${PH} = x$
        
        $\because {AC} = \sqrt{3}r,{PC} = {3r}$ ,则以 ${AB}$ 为底面直径的圆锥容积为 ${V}_{\text{圆锥 }} = \dfrac{1}{3}\pi  \cdot  A{C}^{2} \cdot  {PC}$
        $= \dfrac{1}{3}\pi {\left( \sqrt{3}r\right) }^{2} \cdot  {3r} = {3\pi }{r}^{3}$ ,球取出后水面下降到 ${EF}$ ,水体积为
        ${V}_{\text{水 }} = \dfrac{1}{3}\pi  \cdot  E{H}^{2} \cdot  {PH} = \dfrac{1}{3}\pi {\left( PH\tan {30}^{ \circ  }\right) }^{2}{PH} = \dfrac{1}{9}\pi {x}^{3}.$
        又 ${V}_{\text{水 }} = {V}_{\text{圆锥 }} - {V}_{\text{球 }}$ ,则 $\dfrac{1}{9}\pi {x}^{3} = {3\pi }{r}^{3} - \dfrac{4}{3}\pi {r}^{3}$ ,解得 $x = \sqrt[3]{15}r$ .
    \end{jieda}
    
    \item 鳖臑是我国古代对四个面均为直角三角形的三棱锥的称呼. 如图, 三棱锥 $A - {BCD}$ 是一鳖臑,其中 ${AB} \bot  {BC},{AB} \bot  {BD},{BC} \bot  {CD},{AC} \bot  {CD}$ ,且高 ${AB} = 3\sqrt{2},{BC} = \sqrt{2}, {CD} = \sqrt{6}$ .
    \begin{enumerate*}
        \item 求三棱锥 $A - {BCD}$ 的体积和表面积
        \item 求三棱锥 $A - {BCD}$ 外接球和内切球
    \end{enumerate*}
    \begin{center}
        \includegraphics[width=0.25\linewidth]{bodyxb//src/qiu3.png}
    \end{center}
    \begin{jieda}
        \begin{enumerate}
            \item 由题设可得 ${CD} = \sqrt{3}$ ,而三棱锥的高为 ${AB} = 3\sqrt{2}$ ,
            三棱锥 $A - {BCD}$ 的体积 ${V}_{A - {BCD}} = \dfrac{1}{3} \cdot  {S}_{\bigtriangleup {DBC}} \cdot  {AB} = \dfrac{1}{3} \times  \dfrac{1}{2} \times  \sqrt{6} \times  \sqrt{3} \times  3\sqrt{2} = 3$ ,
            又 ${AC} = \sqrt{{18} + 6} = 2\sqrt{6}$ ,
            三棱锥 $A - {BCD}$ 的表面积 $S = {S}_{\bigtriangleup {BCD}} + {S}_{\bigtriangleup {ABC}} + {S}_{\bigtriangleup {ACD}} + {S}_{\bigtriangleup {ABD}}$
            $= \dfrac{3\sqrt{2}}{2} + 3\sqrt{3} + 3\sqrt{2} + \dfrac{9\sqrt{2}}{2} = 9\sqrt{2} + 3\sqrt{3}.$
            \item 由条件知,可将三棱锥 $A - {BCD}$ 补成一个长方体,则三棱锥的四个顶点也为长方体的
            顶点, 因此长方体的外接球也为三棱锥的外接球. 即为三棱锥外接球的直径.
        \end{enumerate}
    \end{jieda}
        
    
    \item 如图所示，在棱长为 1 的正方体内有两个球相外切且又分别与正方体内切.
    \begin{enumerate*}
        \item 求两球半径之和
        \item 球的半径为多少时，两球体积之和最小
    \end{enumerate*}
    
    \begin{center}
    \includegraphics[width=0.2\textwidth]{bodyxb/src/bo_d5gv3rn7aajc7382nmmg_1_1107_1276_301_337_0.jpg}
    \end{center}
    \begin{jieda}
        \begin{enumerate}
            \item 如图,球心 ${O}_{1}$ 和 ${O}_{2}$ 在 ${AC}$ 上,过 ${O}_{1},{O}_{2}$ 分别作 ${AD},{BC}$ 的垂线交于 $E,F$ .
            
            则由 ${AB} = 1,{AC} = \sqrt{3}$ 得 $A{O}_{1} = \sqrt{3}r,C{O}_{2} = \sqrt{3}R.\therefore r + R + \sqrt{3}\left( {r + R}\right)  = \sqrt{3},\therefore R +r = \dfrac{\sqrt{3}}{\sqrt{3} + 1} = \dfrac{3 - \sqrt{3}}{2}.$
            \begin{center}
            \includegraphics[width=0.25\linewidth]{bodyxb//src/qiu4.png}
        \end{center}
        \item 设两球体积之和为 $V$ ,则 $V = \dfrac{4}{3}\pi \left( {{R}^{3} + {r}^{3}}\right)  = \dfrac{4}{3}\pi \left( {r + R}\right) \left( {{R}^{2} - {Rr} + {r}^{2}}\right) = \dfrac{4}{3}\pi \dfrac{3 - \sqrt{3}}{2}\left\lbrack  {{\left( R + r\right) }^{2} - {3rR}}\right\rbrack   = \dfrac{4}{3}\pi \dfrac{3 - \sqrt{3}}{2}\left\lbrack  {{\left( \dfrac{3\sqrt{3}}{2}\right) }^{2} - {3R}\left( {\dfrac{3\sqrt{3}}{2} - R}\right) }\right\rbrack$
        $= \dfrac{4}{3}\pi \dfrac{3 - \sqrt{3}}{2}\left\lbrack  {3{R}^{2} - \dfrac{3\left( {3 - \sqrt{3}}\right) }{2}R + {\left( \dfrac{3 - \sqrt{3}}{2}\right) }^{2}}\right\rbrack$
        此 $R = \dfrac{3 - \sqrt{3}}{4}$ 时, $V$ 有最小值.
        $\therefore$ 当 $R = r = \dfrac{3 - \sqrt{3}}{4}$ 时,体积之和有最小值.
        \end{enumerate}
    \end{jieda}
\end{enumerate}

表面积与体积练习
\begin{enumerate}
    \item 若正三棱锥的底面边长为 3，侧棱长为 $2\sqrt{3}$ ，则这个正三棱锥的体积是\blkda{$\dfrac{9\sqrt{3}}{4}$}.
    
    \item 在母线长为 1 的圆锥体积最大时，其侧面展开图圆心角 $\varphi$ 等于\blkda{$\dfrac{2\sqrt{6}\pi}{3}$}.

    \item 已知正三棱柱 ${ABC} - {A}_{1}{B}_{1}{C}_{1}$ 的侧面对角线 ${A}_{1}B$ 与侧面 ${AC}{C}_{1}{A}_{1}$ 所成的角为 ${45}^{ \circ  }$ ， ${AB} = {4cm}$ ，则这个棱柱的侧面积是\blkda{$24\sqrt{2}cm^2$}.

    \item 正三棱柱 $P - {ABC}$ 的高为 ${4cm}$ ，侧棱和底面所成的角为 ${45}^{ \circ  }$ ，这个正三棱锥的表面积是\blkda[3]{$\left( {{12}\sqrt{15} + {12}\sqrt{3}}\right) {{{cm}}}^{2}$}.

    \item 正四棱锥底面外接圆半径为 ${10}{{cm}}$ ，斜高为 ${12}{{cm}}$ ，下面数据正确的是 \dotfill{(\kh{D})}
    \xx{高 $2\sqrt{11}{{cm}}$}{侧棱长 $12cm$}{侧面积 ${60}\sqrt{2}{{{cm}}}^{3}$}{对角面面积 ${10}\sqrt{94}{{{cm}}}^{2}$}
    
    \item 若一个圆锥的侧面积是其底面积的 2 倍，则该圆锥的母线与底面所成的角是\blkda{$\dfrac{\pi}{3}$}.

    \item 若一圆锥的侧面展开图为半圆，且面积为 S，则圆锥的底面面积是\blkda{$\dfrac{S}{2}$}.
    
    \begin{center}
    \includegraphics[width=0.2\textwidth]{bodyxb/src/bo_d5gv3rn7aajc7382nmmg_2_1228_992_204_202_0.jpg}
    \end{center}

    \item 由 $18$ 个边长为 $1 cm$ 的小正方体拼成的几何体，几何体的表面积是\blkda{$48cm^2$}.
    
    \item 若两个平行于圆锥底面的平面将圆锥的高分成相等的三段，那么圆锥被分成的三部分的体积的比是\blkda{$1:7:19$}.

    \item 一个三棱锥形容器 $S - {ABC}$ 三条侧棱上各有一个小洞 $D,E,F$ ，且 ${SD} : {DA} = {SE} : {EB} = {CF} : {FS} = 2 : 1$ ，若用它盛水，则最多可盛容器体积的\blkda{$\dfrac{23}{27}$}.

    \item 已知三棱柱 ${ABC} - {A}_{1}{B}_{1}{C}_{1}$ 中，底面 ${ABC}$ 为等边三角形， $A{A}_{1}$ 与 ${AB},{AC}$ 的夹角均为 ${60}^{ \circ  }$ ， ${AB} = 4$ ， ${AA}_{1} = 6$ ，则此棱柱的体积是\blkda{$24\sqrt{2}$}.

    \item 如果一个圆锥和一个半球有公共底面，圆锥的体积恰好等于半球的体积，那么这个圆锥的轴截面的顶角的余弦值是\blkda{$\dfrac{3}{5}$}.

    \item 已知圆柱形容器的内壁底面半径为 $5{{cm}}$ ，两个直径为 $5{{cm}}$ 的玻璃小球都浸没于这个容器的水中.若取出这两个小球，则容器内的水面下降\blkda{$\dfrac{5}{3}$} ${cm}$ .
\end{enumerate}